\documentclass[11pt]{article}
\usepackage[margin=1in]{geometry}
\usepackage{amsmath,amssymb,amsthm,mathtools,bm,mathrsfs}
\usepackage[hidelinks]{hyperref}
\newtheorem{theorem}{Theorem}[section]
\newtheorem{lemma}[theorem]{Lemma}
\newtheorem{proposition}[theorem]{Proposition}
\newtheorem{corollary}[theorem]{Corollary}
\newcommand{\SU}{\operatorname{SU}}
\newcommand{\tr}{\operatorname{tr}}
\newcommand{\Var}{\operatorname{Var}}
\newcommand{\Rea}{\operatorname{Re}}
\newcommand{\dd}{\,\mathrm d}
\newcommand{\HH}{\mathscr H}
\newcommand{\KK}{\mathscr K}
\newcommand{\LL}{\mathscr L}
\newcommand{\VV}{\mathscr V}
\newcommand{\EE}{\mathsf E}
\newcommand{\PP}{\mathsf P}
\newcommand{\NN}{\mathcal N}
\title{Extensive true-vacuum quantum blocking in the full spatial Wilson box\\
Exact path-product maps, associative memory, and physical spectral channels}
\author{}
\date{8 September 2026}
\begin{document}
\maketitle

\section{Original objects and the finite interacting operator}

Let $L\geq2$ be an integer. The vertex set is
$\mathsf V_L=\{-L,-L+1,\ldots,L\}^3$. For each $i\in\{1,2,3\}$,
the physical edge $e=(n,i)$ goes from $n$ to $n+\mathbf e_i$ when both
vertices are present. Its variable is $U_i(n)\in\SU(2)$. Inverse traversal
uses $U_i(n)^{-1}$. Every such edge belongs to $\EE_L$. Every elementary
face $p=(n;i,j)$, $i<j$, with four vertices in the box belongs to $\PP_L$;
its ordered holonomy and trace are
\begin{equation}\label{eq:plaquette}
 U_p=U_i(n)U_j(n+\mathbf e_i)U_i(n+\mathbf e_j)^{-1}U_j(n)^{-1},
 \qquad W_p=\tr U_p.
\end{equation}
The counts are $N=3(2L)(2L+1)^2$ and $M=3(2L)^2(2L+1)$: fix a
direction or a face plane and count its initial and transverse coordinates.
The configuration space is $\mathcal Q=\SU(2)^{\EE_L}$ with product
Haar probability measure $\lambda$. Put
\begin{gather}
 T_a=-i\sigma_a/2,\qquad
 X_{e,a}f=\left.\frac{d}{dt}f(\ldots,e^{tT_a}U_e,\ldots)\right|_{t=0},
 \qquad E_e=-\sum_{a=1}^3X_{e,a}^2,\label{eq:fields}\\
 H=\kappa\sum_{e\in\EE_L}E_e+b\sum_{p\in\PP_L}(2-W_p),\qquad
 \kappa=\frac{2g^2}{a},\quad b=\frac{1}{2g^2a},\quad
 \xi=\frac b\kappa=\frac1{4g^4},\quad a,g>0.\label{eq:H}
\end{gather}
Here $a$ in $a,g>0$ denotes spatial spacing; the color index in
\eqref{eq:fields} is summed only when indicated. Haar integration by parts
makes $X_{e,a}$ skew-adjoint. The Pauli identities give
$-2\tr(T_aT_b)=\delta_{ab}$ and $-\sum_aT_a^2=3I/4$.
The original metric $c(A,B)=-\tr(AB)/2$ has orthonormal basis $2T_a$;
hence $X^c=2X$, $\Delta^c=4\Delta^T$, and the same kinetic term is
$-(\kappa/4)\sum_e\Delta_e^c$. All faces, including boundary faces
and faces in every plane, remain in \eqref{eq:H}.
In particular the exact scalar term $2bM=M/(g^2a)$ is part of $H$.

Gauge transformations are $(h\cdot U)_e=h_{s(e)}^{-1}U_eh_{t(e)}$.
The form of $H$ on $H^1(\mathcal Q)$ is
\[
 q_H(u,v)=\kappa\sum_{e,a}\int\overline{X_{e,a}u}X_{e,a}v\dd\lambda
        +b\sum_p\int(2-W_p)\overline uv\dd\lambda.
\]
It is closed, since the smooth potential is bounded between $0$ and
$4bM$. Its associated self-adjoint operator has domain $H^2$ by elliptic
regularity; smooth functions are operator and form cores. Compact Sobolev
embedding gives compact resolvent. Taking the modulus of a minimizing
vector does not increase its form: apply the chain rule first to
$(|u|^2+\epsilon^2)^{1/2}$ and then let $\epsilon$ decrease to zero.
The resulting nonnegative lowest eigenfunction is smooth by elliptic
regularity and strictly positive by the elliptic strong maximum principle.
Write it as $\psi>0$, with $H\psi=E_0\psi$ and $\int\psi^2\dd\lambda=1$.
This fixes the state norm; none of the physical coefficients is changed.

Set $\rho=\psi^2$, $\mu=\rho\lambda$, and $\HH=L^2(\mu)$.
Multiplication $\mathcal Uu=\psi u$ is a unitary $\HH\to L^2(\lambda)$,
with inverse division by $\psi$. The product rule and $H\psi=E_0\psi$ give
\begin{align}
 \LL&=\mathcal U^{-1}(H-E_0)\mathcal U
 =-\kappa\sum_{e,a}\bigl[X_{e,a}^2+(X_{e,a}\log\rho)X_{e,a}\bigr],
 \label{eq:gs}\\
 q(u,v)&=\kappa\sum_{e,a}\int\overline{X_{e,a}u}X_{e,a}v\dd\mu,
 \qquad D(q)=\VV=H^1(\mathcal Q).\label{eq:q}
\end{align}
Integration by parts proves the form identity. Multiplication and division
by the smooth positive $\psi$ preserve every integer Sobolev space, so
the form and operator domains are respectively $H^1,H^2$.
If $q(u,u)=0$, all its derivatives vanish and connectedness makes $u$
constant. This proves uniqueness of $\psi$, and its gauge invariance
follows because gauge transformations preserve $H$, positivity and norm.
It also proves that the finite first nonzero eigenvalue $\delta$ of
$\LL$ is strictly positive. No uniform lower bound for $\delta$ is used
in the uniform estimates below. Inner products are conjugate-linear in
their first arguments.

\section{A uniform single-link estimate for the actual vacuum}

This section supplies estimates independent of $L$, with every $b>0$
allowed. Define the constants
\begin{equation}\label{eq:uniformconstants}
 D=16b,\quad d=D/\kappa=16\xi=4/g^4,\quad
 t_0=\frac{2\log12}{\kappa},\quad
 R=3e^{Dt_0}=3\,12^{2d},\quad \alpha=R^{-2},
\end{equation}
and
\begin{equation}\label{eq:G}
 G=R\exp\!\left(\frac d{1+d}\right)
       \left(1+\frac{2d}{1+d}\right)\sqrt{\frac{1+d}{2}},
 \qquad C_*=\frac32+2G.
\end{equation}
Here $D$ has units of energy and $t_0$ of inverse energy.
The constants $d,R,\alpha,G,C_*$ are dimensionless and depend on the
original coupling $g$, but not on the number of links, the blocking
length, or $a$ separately.

\begin{lemma}\label{lem:singleheat}
For the single-link heat operator $P_t=e^{-t\kappa E}$, its kernel at
$t=t_0$ lies between $1/2$ and $3/2$. For every real smooth $f$,
\begin{equation}\label{eq:reverse}
 2\kappa t\sum_a|X_aP_tf|^2
 \leq P_t(f^2)-(P_tf)^2,\qquad
 \|\nabla P_tf\|_\infty\leq\frac{\|f\|_\infty}{\sqrt{2\kappa t}}.
\end{equation}
\end{lemma}
\begin{proof}
The spin-$j$ representation has dimension $2j+1$ and Casimir $j(j+1)$
in \eqref{eq:fields}. One can verify the latter by the eigenvalues
$m=-j,\ldots,j$ of $iT_3$ and the raising/lowering coefficients
$\sqrt{(j\mp m)(j\pm m+1)}$; their quadratic sum is $j(j+1)$.
Peter--Weyl completeness therefore gives the heat kernel
$p_t(U)=\sum_{j\geq0}(2j+1)e^{-\kappa tj(j+1)}\chi_j(U)$.
This series is absolutely uniformly convergent for $t>0$, since
$|\chi_j|\leq2j+1$. Writing $k=2j$, using
$j(j+1)=k(k+2)/4\geq k/2$ and $(k+1)^2\leq4^k$ for $k\geq1$, gives
\[
 |p_{t_0}-1|\leq\sum_{k\geq1}4^k e^{-\kappa t_0k/2}
 =\sum_{k\geq1}3^{-k}=\frac12.
\]
For the gradient assertion, the bi-invariant Casimir commutes with
each $X_a$: the commutator of its sum of squares with $X_b$ contracts
antisymmetric Lie structure constants with $X_aX_c+X_cX_a$ and is zero.
Thus $X_aP_s=P_sX_a$ on smooth functions, and Jensen's inequality gives
$\Gamma(P_sf)\leq P_s\Gamma(f)$ for $\Gamma(f)=\sum_a|X_af|^2$.
Differentiate $P_s[(P_{t-s}f)^2]$ and use the product rule:
\[
 P_t(f^2)-(P_tf)^2=2\kappa\int_0^tP_s\Gamma(P_{t-s}f)\dd s
 \geq2\kappa t\Gamma(P_tf).
\]
The last inequality applies the preceding gradient contraction with
input $P_{t-s}f$. The supremum estimate follows because
$P_t(f^2)\leq\|f\|_\infty^2$. This is the exact single-group version
of the reverse Poincare inequality \cite[Eq.~(1.7)]{BL}, with the time
and generator coefficient displayed explicitly.
\end{proof}

\begin{theorem}\label{thm:localvacuum}
Fix one physical link $e$ and write all other variables as $y$. For all
$u,u'\in\SU(2)$ and $y$,
\begin{equation}\label{eq:ratio}
 R^{-1}\leq\frac{\psi(u,y)}{\psi(u',y)}\leq R,\qquad
 \alpha\leq\frac{\rho(u,y)}{\int\rho(v,y)\dd v}\leq\alpha^{-1}.
\end{equation}
Moreover,
\begin{equation}\label{eq:scorebound}
 \left(\sum_a|X_{e,a}\log\psi|^2\right)^{1/2}\leq G,
 \qquad |\nabla_e\log\rho|\leq2G.
\end{equation}
All bounds are uniform in $e$ and $L$ at the fixed positive $a,g$.
\end{theorem}
\begin{proof}
Let $V_e=b\sum_{p\ni e}(2-W_p)$ and let $H_{\neg e}$ contain all
other kinetic terms and all faces not incident to $e$. There are at
most four incident faces in the full three-dimensional box, including
the boundary cases with fewer faces. Therefore $0\leq V_e\leq D$,
and the exact operator is
\[
 H=A_e+V_e,\qquad A_e=\kappa E_e+H_{\neg e},\qquad
 T_t=e^{-tA_e}=P_t^{(e)}\otimes e^{-tH_{\neg e}}.
\]
The two factors of $T_t$ commute because $H_{\neg e}$ is independent
of $u$. They have positive kernels. The parabolic comparison principle
with the bounded multiplication $0\leq V_e\leq D$ gives, for $f\geq0$,
\begin{equation}\label{eq:comparison}
 e^{-Dt}T_tf\leq e^{-tH}f\leq T_tf.
\end{equation}
For example $T_tf$ is a supersolution for
$\partial_t+A_e+V_e$, while $e^{-Dt}T_tf$ is a subsolution; their
initial data agree. The maximum principle gives \eqref{eq:comparison}.
This proves the comparison for this compact group product directly;
no Euclidean configuration-space identification is needed.

Apply it to $\psi$. At $t_0$, positivity of the other-link kernel and
Lemma~\ref{lem:singleheat} imply
\[
 \frac{(T_{t_0}\psi)(u,y)}{(T_{t_0}\psi)(u',y)}\leq3.
\]
Since $e^{-tH}\psi=e^{-tE_0}\psi$, comparison at both $u,u'$ gives
$\psi(u,y)/\psi(u',y)\leq3e^{Dt_0}=R$.
Interchanging $u,u'$ proves both sides. Squaring and averaging the
denominator over $u'$ proves the conditional density bounds.
The extensive energy $E_0$ cancels from this ratio exactly.

Here is a derivative proof that does not differentiate an unknown
vacuum estimate. If a positive function $F(u,y)$ satisfies the first
ratio bound, so does
$f_t(u)=(e^{-tH_{\neg e}}F(u,\cdot))(y)$ at each fixed $y$, by
positivity. Hence $\|f_t\|_\infty\leq R\min f_t\leq RP_t f_t(u)$.
Lemma~\ref{lem:singleheat} yields
\[
 |\nabla_e T_t\psi|\leq\frac{R}{\sqrt{2\kappa t}}T_t\psi.
\]
The function $(e^{-tH_{\neg e}}(V_e\psi))(u,y)$ is nonnegative and
bounded above by $D(e^{-tH_{\neg e}}\psi)(u,y)$. Its supremum in $u$
is consequently at most $DR\min_u(e^{-tH_{\neg e}}\psi)(u,y)$.
The same heat estimate gives
\[
 |\nabla_e T_t(V_e\psi)|\leq\frac{DR}{\sqrt{2\kappa t}}T_t\psi.
\]
The exact Duhamel identity on the smooth eigenvector is
\begin{equation}\label{eq:duhamel}
 \psi=e^{E_0t}T_t\psi-\int_0^t e^{E_0\tau}T_\tau(V_e\psi)\dd\tau.
\end{equation}
It follows either by differentiating $e^{E_0t}T_t\psi$ and using
$A_e\psi=(E_0-V_e)\psi$, or by the bounded-potential Duhamel formula.
Comparison gives $e^{E_0\tau}T_\tau\psi\leq e^{D\tau}\psi$.
The two gradient estimates, whose $\tau^{-1/2}$ singularity is
integrable, justify differentiation of \eqref{eq:duhamel} and give
\[
 \frac{|\nabla_e\psi|}{\psi}
 \leq R\left[\frac{e^{Dt}}{\sqrt{2\kappa t}}
       +D\int_0^t\frac{e^{D\tau}}{\sqrt{2\kappa\tau}}\dd\tau\right]
 \leq\frac{Re^{Dt}(1+2Dt)}{\sqrt{2\kappa t}}.
\]
Choose the actual time $t=(\kappa+D)^{-1}$. This last expression is
exactly $G$ in \eqref{eq:G}. Finally $X\log\rho=2X\log\psi$.
No gap bound, perturbation disk, or correlation factorization was used.
\end{proof}

\section{A specified extensive spatial blocking sequence}

Let $m\geq1$ divide $2L$, and put $q_m=2L/m$. Retain as skeleton links
those direction-$i$ fine edges whose two transverse coordinates lie in
$-L+m\{0,1,\ldots,q_m\}$. Every longitudinal fine edge on each such
line is included. Denote this subset by $R_m$ and its complement by $S_m$.
Its exact count is
\begin{equation}\label{eq:counts}
 |R_m|=3(2L)(q_m+1)^2,\qquad
 |S_m|=3(2L)\bigl[(2L+1)^2-(q_m+1)^2\bigr].
\end{equation}
Thus a fixed $m>1$ eliminates an extensive number of links as the box
increases. If $m\mid n\mid2L$, then $R_n\subset R_m$. In particular
$m=1,2,4,\ldots,2^k$ gives a repeated blocking sequence whenever $2^k\mid2L$.
The Hamiltonian still has all $N$ link and $M$ face terms.

The coarse vertices are the points
$x\in[-L,L]^3$ whose coordinates all lie in $-L+m\mathbb Z$.
A coarse direction-$i$ edge $\gamma=(x,i)$ is the ordered path of $m$
fine edges from $x$ to $x+m\mathbf e_i$. The number of coarse edges is
\[
 N_m=3q_m(q_m+1)^2,\qquad |R_m|=mN_m.
\]
The chains have disjoint physical-edge supports and meet only at coarse
vertices. Write their variables as $u_{\gamma,1},\ldots,u_{\gamma,m}$
in positive order, and define
\begin{equation}\label{eq:path}
 z_\gamma=u_{\gamma,1}\cdots u_{\gamma,m},\qquad
 \pi_m:\mathcal Q\longrightarrow\mathcal Q_m:=\SU(2)^{N_m}.
\end{equation}
No original coordinate was replaced by a unit coarse spacing: the physical
coarse spacing is $ma$, while the original coordinates remain $x$.

For every chain put $v_{\gamma,k}=u_{\gamma,k}$ for $k<m$.
Let $y$ comprise all $v$ and all untouched variables on $S_m$.
The full typed coordinate map is
\[
 \Phi_m:\mathcal Q\longrightarrow\mathcal Y_m\times\mathcal Q_m,
 \quad U\longmapsto(y,z),\qquad
 \mathcal Y_m=\SU(2)^{S_m}\times
                \prod_{\gamma=1}^{N_m}\SU(2)^{m-1}.
\]
It is a global smooth bijection. Its inverse
$\Phi_m^{-1}:\mathcal Y_m\times\mathcal Q_m\to\mathcal Q$ uses
\[
 u_{\gamma,m}=(v_{\gamma,1}\cdots v_{\gamma,m-1})^{-1}z_\gamma,
\]
and keeps every other variable. Substituting the original ordered
product cancels its prefix on the left and recovers $u_{\gamma,m}$.
Conversely multiplying the retained prefix by the displayed last link
recovers $z_\gamma$. Every $v$ and every $S_m$ coordinate is fixed,
so both compositions are the identity on their respective full spaces.
There are $|S_m|+(m-1)N_m=N-N_m$ factors in $\mathcal Y_m$;
none of those coordinates is omitted from $\Phi_m$.
Integrating $u_{\gamma,m}$ last,
left invariance of Haar proves the exact measure identity
\begin{equation}\label{eq:haarcoordinates}
 \dd\lambda(U)=\dd\lambda_y(y)\dd\lambda_m(z).
\end{equation}
This proves the map and its measure factor, rather than postulating
independent coarse links in the interacting vacuum.

The complete kinetic operator in these coordinates is as follows.
Define the prefix and the real
orthogonal matrix $O_{\gamma,k}$ by
\[
 p_{\gamma,k}=v_{\gamma,1}\cdots v_{\gamma,k-1},\qquad
 p_{\gamma,k}T_a p_{\gamma,k}^{-1}
   =\sum_b O_{\gamma,k;ba}T_b.
\]
Orthogonality follows by applying $-2\tr(AB)$ to both conjugated factors;
its determinant is $1$ by connectedness. Let $D_{\gamma,b}$ be left
multiplication differentiation on $z_\gamma$, with the same $T_b$.
For $k<m$ the transformed fine vector field is
\[
 X_{\gamma,k,a}=X_{v_{\gamma,k},a}
            +\sum_bO_{\gamma,k;ba}D_{\gamma,b};
\]
for $k=m$ it is $\sum_bO_{\gamma,m;ba}D_{\gamma,b}$. The prefix
$p_{\gamma,k}$ does not involve $v_{\gamma,k}$, and no $O$ depends on $z$.
The full sum of squares is consequently
\begin{align}\label{eq:fulltensor}
 \sum_{e,a}X_{e,a}^2
 ={}&\sum_{e\in S_m,a}X_{e,a}^2
 +\sum_{\gamma}\left[\sum_{k<m,a}X_{v_{\gamma,k},a}^2
                      +m\sum_bD_{\gamma,b}^2\right]\nonumber\\
 &+2\sum_{\gamma,k<m,a,b}O_{\gamma,k;ba}
                    X_{v_{\gamma,k},a}D_{\gamma,b}.
\end{align}
Each square before expansion is the square of the displayed original
field, so this formula includes all mixed derivatives with their signs
and coefficients. Every original $W_p$ in \eqref{eq:H} is evaluated at
the explicit inverse coordinate map. It is not replaced by a coarse face
trace or truncated to selected words.

Gauge transformations on a chain cancel at its internal vertices and give
$z_\gamma\mapsto h_x^{-1}z_\gamma h_{x+m\mathbf e_i}$.
The marginal density on $R_m$ is invariant under the internal gauge
transformations, by gauge invariance of $\rho$ and Haar substitution on
$S_m$. A non-coarse skeleton vertex is bivalent in the skeleton. For a
fixed chain with endpoints fixed, set successively
$h_{x+k\mathbf e_i}=(u_1\cdots u_k)^{-1}$, $k<m$.
This sends its first $m-1$ links to the identity and its last to $z_\gamma$.
Different chain interiors are disjoint, so these choices are simultaneous.
It follows that this skeleton marginal is a function of $z$ alone.
This statement concerns the marginal; the full conditional distribution
of $y$ given $z$ need not factorize.

\section{The actual coarse density, domains, and conditional coupling}

Let $\widetilde\rho_m(y,z)=\rho(U(y,z))$ and define
\begin{gather}
 r_m(z)=\int\widetilde\rho_m(y,z)\dd\lambda_y,
 \quad \mu_m=r_m\lambda_m,\quad
 p_m(y\mid z)=\widetilde\rho_m(y,z)/r_m(z),\label{eq:coarsedensity}\\
 \KK_m=L^2(\mu_m),\quad J_m f=f\circ\pi_m,\quad
 P_m=J_mJ_m^*,\quad Q_m=I-P_m,\nonumber\\
 (J_m^*u)(z)=\int p_m(y\mid z)u(U(y,z))\dd\lambda_y.\label{eq:J}
\end{gather}
In particular the types are $J_m:\KK_m\to\HH$ and
$J_m^*:\HH\to\KK_m$. Their full coordinate expressions are
$J_mf(U)=f(z(U))$ and the fibre integral \eqref{eq:J};
the operator $P_m:\HH\to J_m\KK_m\subset\HH$ is their composition.
Smoothness and strict positivity follow from compact integration.
Fubini proves the adjoint formula and $J_m^*J_m=I$; Jensen proves that
$J_m^*$ is a contraction. Thus $P_m$ is the true-vacuum orthogonal
projection. All these maps intertwine the stated fine and coarse gauge
actions. The restriction of the conditional projection onto skeleton
functions, on the fine invariant Hilbert space, equals this path projection:
its output is internally gauge invariant, hence depends only on $z$, as
proved in the preceding section. The exact invariant Hilbert-space
unitary, its almost-everywhere inverse and the equality of the full
complementary operator domains are proved in
Proposition~\ref{prop:physicalbridge}.

For each chain and color set, in the original coordinates,
\begin{equation}\label{eq:Y}
 Y_{\gamma,b}=\frac1m\sum_{k=1}^m\sum_a
                   O_{\gamma,k;ba}X_{\gamma,k,a}.
\end{equation}
Then $Y_{\gamma,b}J_mf=J_mD_{\gamma,b}f$.
Indeed each summand differentiates $z_\gamma$ in the direction
$\operatorname{Ad}_{p_{\gamma,k}}T_a$, and orthogonality of $O$ gives
one copy of $D_{\gamma,b}$ per $k$.
Also $Y_{\gamma,b}$ has zero Haar divergence: every coefficient
$O_{\gamma,k;ba}$ is independent of its own differentiated fine edge.
Consequently integration by parts against a coarse smooth test function
gives the exact derivative identities
\begin{align}
 D_{\gamma,b}\log r_m&=J_m^*(Y_{\gamma,b}\log\rho),\label{eq:marginalscore}\\
 D_{\gamma,b}J_m^*u
   &=J_m^*(Y_{\gamma,b}u)+J_m^*(\beta_{\gamma,b}u),\label{eq:derivativeJ}\\
 \beta_{\gamma,b}&=Y_{\gamma,b}\log\rho
                       -J_mD_{\gamma,b}\log r_m,\qquad J_m^*\beta_{\gamma,b}=0.
 \label{eq:beta}
\end{align}
For example integrate $Y_{\gamma,b}(J_m\varphi\,\rho)$ over the full
configuration space; its integral is zero. The resulting identity is
$\int(D_{\gamma,b}\varphi)r_m=-\int\varphi J_m^*(Y_{\gamma,b}\log\rho)r_m$,
which proves \eqref{eq:marginalscore}. Apply the same calculation to
$u\rho$ to prove \eqref{eq:derivativeJ}. This also avoids differentiating
fibre coordinates while holding an incompatible fine coordinate fixed.

Theorem~\ref{thm:localvacuum}, the average in \eqref{eq:Y}, and conditional
Jensen give
\begin{equation}\label{eq:coarsescorebound}
 |Y_\gamma\log\rho|\leq2G,\qquad
 |D_\gamma\log r_m|\leq2G,\qquad |\beta_\gamma|\leq4G.
\end{equation}
Here each absolute value is the Euclidean norm of its three color
components. The first inequality uses the average of $m$ rotated vectors,
each with norm at most $2G$; it does not assume independence.

\begin{proposition}\label{prop:forms}
The exact lifted form and its generator are
\begin{align}
 a_m(f,g)&=q(J_mf,J_mg)
    =\kappa m\sum_{\gamma,b}\int\overline{D_{\gamma,b}f}D_{\gamma,b}g\dd\mu_m,
       &D(a_m)&=\VV_m=H^1(\mathcal Q_m),\label{eq:am}\\
 A_m&=-\kappa m\,r_m^{-1}\sum_{\gamma,b}D_{\gamma,b}(r_mD_{\gamma,b}),
       &D(A_m)&=H^2(\mathcal Q_m).\label{eq:Am}
\end{align}
The maps $J_m,J_m^*$ preserve all integer Sobolev spaces, and $P_m,Q_m$
preserve $\VV$. On $\VV_m$ the full eliminated coupling is
\begin{equation}\label{eq:Bm}
 B_mf=-\kappa m\sum_{\gamma,b}\beta_{\gamma,b}J_mD_{\gamma,b}f,
 \quad J_m^*B_mf=0,\quad \LL J_mf=J_mA_mf+B_mf\quad(f\in H^2).
\end{equation}
In particular the retained kinetic coefficient is exactly $\kappa m$.
\end{proposition}
\begin{proof}
The fine derivative of $J_mf$ on edge $(\gamma,k)$ is the orthogonal
rotation $O_{\gamma,k}^TD_\gamma f$; on $S_m$ it is zero.
Summing the squared derivatives over the $m$ edges proves \eqref{eq:am}.
The smooth positive weight gives closedness and the $H^2$ elliptic domain
of \eqref{eq:Am}. The compact smooth coordinate bijection
\eqref{eq:haarcoordinates} preserves each integer Sobolev class with
finite norm bounds, since its derivatives and inverse derivatives are
bounded. The fibre integral with smooth positive kernel $p_m$ has bounded
derivatives of every finite order. This proves the Sobolev mapping claims.

There is also a useful explicit first-order estimate. Define
\[
 \eta_m=\kappa m\left\|\sum_{\gamma,b}\beta_{\gamma,b}^2\right\|_\infty
             \leq16\kappa mN_mG^2.
\]
Orthogonality of the rotations and Cauchy--Schwarz in the $k$ sum give
the pointwise inequality
\[
 m\sum_{\gamma,b}|Y_{\gamma,b}u|^2
 \leq\sum_{e\in R_m,a}|X_{e,a}u|^2.
\]
Conditional Jensen and \eqref{eq:derivativeJ} consequently give
\begin{equation}\label{eq:projectionbound}
 a_m(J_m^*u,J_m^*u)\leq2q(u,u)+2\eta_m\|u\|_\mu^2.
\end{equation}
This extends by smooth approximation. The same derivative identities
show that $J_mf\in H^1$ implies $f\in H^1$, so the retained form space
has exactly the claimed domain.
On a lifted smooth $f$, \eqref{eq:fulltensor} gives
$\sum_{e,a}X_{e,a}^2J_mf=mJ_m\sum_{\gamma,b}D_{\gamma,b}^2f$.
The drift in \eqref{eq:gs} is
$-\kappa m\sum_{\gamma,b}(Y_{\gamma,b}\log\rho)J_mD_{\gamma,b}f$.
Subtract \eqref{eq:Am}, and use \eqref{eq:beta}; this proves
\eqref{eq:Bm} for smooth $f$. Density gives the stated $H^2$ identity
and the bounded map $B_m:\VV_m\to Q_m\HH$, with
\begin{equation}\label{eq:Bbound}
 \|B_mf\|_\mu^2\leq\eta_m a_m(f,f).
\end{equation}
\end{proof}

The full correlation tensor of the induced first-order terms is
\begin{equation}\label{eq:Fisher}
 I_{\gamma b,\eta c}(z)=J_m^*(\beta_{\gamma,b}\beta_{\eta,c})(z),\quad
 \|B_mf\|^2=(\kappa m)^2\int
 \sum_{\gamma,b,\eta,c}I_{\gamma b,\eta c}
              \overline{D_{\gamma,b}f}D_{\eta,c}f\dd\mu_m.
\end{equation}
Every off-diagonal index is retained. This is the conditional covariance
of the vectors $Y_{\gamma,b}\log\rho$. It is positive semidefinite
because its quadratic form is a conditional expectation of a squared
absolute value. Positivity does not assert that its different entries
vanish or that it is uniformly diagonal-dominant.

There is an exact Haar-space description of the instantaneous marginal
operator. Multiplication $f\mapsto\sqrt{r_m}f$ is a unitary
$\KK_m\to L^2(\lambda_m)$, with inverse
$v\mapsto v/\sqrt{r_m}$. Both maps preserve $H^1$ and $H^2$,
since $r_m$ is smooth and strictly positive on the compact group
product. Direct product differentiation gives
\begin{equation}\label{eq:marginalpotential}
 \sqrt{r_m}A_m r_m^{-1/2}
 =-\kappa m\sum_{\gamma,b}D_{\gamma,b}^2
       +\kappa m\frac{\sum_{\gamma,b}D_{\gamma,b}^2\sqrt{r_m}}{\sqrt{r_m}}.
\end{equation}
The last scalar function equals
$\kappa m\sum_{\gamma,b}[D_{\gamma,b}^2\log r_m/2
 +(D_{\gamma,b}\log r_m)^2/4]$.
It depends on the entire induced density. No equality with a sum of
one-plaquette Wilson terms is asserted. Moreover \eqref{eq:marginalpotential}
is only the instantaneous compression; the next section gives the rest
of the physical operator relation.

\section{Exact memory, resolvent reconstruction, and physical norms}

Set $\HH_{Q_m}=Q_m\HH$ and $\VV_{Q_m}=\VV\cap\HH_{Q_m}$.
The restriction $c_m(h,k)=q(h,k)$ on $\VV_{Q_m}$ is closed and densely
defined. Closedness uses the $L^2$-closed subspace, and density follows
by applying $Q_m$ to smooth approximants, using \eqref{eq:projectionbound}.
Define its self-adjoint operator $C_m$ by
\begin{align}\label{eq:Cdomain}
 D(C_m)=\{h\in\VV_{Q_m}:&\ \exists v\in\HH_{Q_m}\ \text{such that}
 \ c_m(k,h)=\langle k,v\rangle\ \forall k\in\VV_{Q_m}\},
 \qquad C_mh=v.
\end{align}
For smooth $h\in\HH_{Q_m}$ this is $Q_m\LL h$.
Every such $h$ is orthogonal to $1=J_m1$, so $C_m\geq\delta I>0$.
For form-domain arguments no formal compression domain is substituted
for \eqref{eq:Cdomain}.

For $f,g\in\VV_m$ and $h,k\in\VV_{Q_m}$, integration by parts on a
smooth core, followed by continuity, gives
\begin{equation}\label{eq:block}
 q(J_mf+h,J_mg+k)=a_m(f,g)+\langle B_mf,k\rangle
                         +\langle h,B_mg\rangle+c_m(h,k).
\end{equation}
For smooth $h\in\HH_{Q_m}$ the adjoint coupling is
\begin{equation}\label{eq:Bstar}
 B_m^*h=\kappa m\,r_m^{-1}\sum_{\gamma,b}
           D_{\gamma,b}\bigl(r_mJ_m^*(\beta_{\gamma,b}h)\bigr).
\end{equation}
To prove the sign, integrate the negative derivative in \eqref{eq:Bm}
against $h$ using Haar integration by parts in $z$. Differentiating the
conditional integral by \eqref{eq:derivativeJ} expands each derivative
in \eqref{eq:Bstar} into terms with $Yh$, $Y\beta$, $\beta^2h$ and
$(D\log r_m)\beta h$. All coefficients are smooth and bounded in a fixed
box, so $B_m^*:\VV_{Q_m}\to\KK_m$ is bounded. Without that regularity
the adjoint is the bounded dual map $\HH_{Q_m}\to\VV_m^*$.

For $z>0$ define on $\VV_m$ the sesquilinear form
\begin{equation}\label{eq:schur}
 s_{m,z}(f,g)=a_m(f,g)+z\langle f,g\rangle
             -\langle B_mf,(C_m+z)^{-1}B_mg\rangle.
\end{equation}
The use of $z$ as a positive spectral shift here is distinct from the
tuple $z_\gamma$ of coarse link variables. The argument of a resolvent
always denotes the scalar shift.

\begin{theorem}\label{thm:schur}
The form $s_{m,z}$ is closed with domain $\VV_m$, and its associated
operator $S_{m,z}$ satisfies $S_{m,z}\geq zI$. Its exact full resolvent
and reconstruction maps are
\begin{align}
 J_m^*(\LL+z)^{-1}J_m&=S_{m,z}^{-1},\label{eq:resolvent}\\
 (\LL+z)^{-1}J_mg
 &=\bigl[J_m-(C_m+z)^{-1}B_m\bigr]S_{m,z}^{-1}g.\label{eq:reconstruction}
\end{align}
Every product in these formulas is defined by the displayed form domains.
\end{theorem}
\begin{proof}
Completing the square in \eqref{eq:block} gives
\begin{align}\label{eq:square}
 q(J_mf+h,J_mf+h)+z(\|f\|^2+\|h\|^2)
 =s_{m,z}(f,f)+\|(C_m+z)^{1/2}[h+(C_m+z)^{-1}B_mf]\|^2.
\end{align}
Its unique minimizing $h$ is $-(C_m+z)^{-1}B_mf$, which lies in
$D(C_m)\subset\VV_{Q_m}$. Nonnegativity of $q$ proves
$s_{m,z}(f,f)\geq z\|f\|^2$. Apply \eqref{eq:projectionbound} to
$u=J_mf+h$. It gives
\[
 a_m(f,f)+\kappa\|f\|^2\leq2q(u,u)+(2\eta_m+\kappa)\|u\|^2
 \leq K_{m,z}\,[q(u,u)+z\|u\|^2],
\]
where $K_{m,z}=\max\{2,(2\eta_m+\kappa)/z\}$ in the form norm with
reference energy $\kappa$. This constant is dimensionless.
Taking the minimizing $h$ proves a lower form-norm bound. The upper
bound $s_{m,z}(f,f)\leq a_m(f,f)+z\|f\|^2$ follows by choosing $h=0$.
Their equivalence proves closedness. For $v=(\LL+z)^{-1}J_mg$, its
form-domain decomposition $v=J_mf+h$ is legitimate. Testing the weak
resolvent equation against $k\in\VV_{Q_m}$ gives
$h=-(C_m+z)^{-1}B_mf$; testing against $J_mf'$ gives
$s_{m,z}(f',f)=\langle f',g\rangle$. This proves both identities.
\end{proof}

The notation in \eqref{eq:schur} is a closed-form implementation of the
Feshbach--Schur map. The general method and reconstruction are established
in \cite[Theorem~1.2, Eqs.~(1.10)--(1.16)]{DSS}. That theorem's bounded
effective operator assumption is not silently applied to our infinite-rank
projection; the coercivity and domain proof above supplies the needed
extension for this particular operator.

For smooth initial $f_0$, define
$K_m(t)=J_m^*e^{-t\LL}J_m$, $f(t)=K_m(t)f_0$.
The full evolution and \eqref{eq:block} give
\begin{align}
 f'(t)&=-A_m f(t)+\int_0^t B_m^*e^{-(t-s)C_m}B_m f(s)\dd s,
       &f(0)&=f_0,\label{eq:memory}\\
 K_m(t)&=J_m^*e^{-t\LL}J_m,\qquad
 \int_0^\infty e^{-zt}K_m(t)\dd t=S_{m,z}^{-1}.\label{eq:laplace}
\end{align}
Here is the domain justification and the sign calculation. Smooth initial
data lie in all powers of the compact elliptic $\LL$. Spectral calculus
and elliptic regularity keep $v(t)=e^{-t\LL}J_mf_0$, $f(t)=J_m^*v(t)$
and $h(t)=Q_mv(t)$ smooth, continuously in the required Sobolev norms.
They satisfy $f'=-A_mf-B_m^*h$, $h'=-B_mf-C_mh$, $h(0)=0$.
Variation of constants gives
$h(t)=-\int_0^te^{-(t-s)C_m}B_mf(s)\dd s$.
Since $\sup_{x\geq0}\sqrt{x}e^{-\tau x}=(2e\tau)^{-1/2}$,
$e^{-\tau C_m}$ maps $\HH_{Q_m}$ to its form space with norm at most
$\sqrt\kappa+(2e\tau)^{-1/2}$ when the squared form norm is
$c_m(h,h)+\kappa\|h\|^2$. The singularity is integrable; the bounded form-space
map \eqref{eq:Bstar} makes \eqref{eq:memory} a Bochner integral in
$\KK_m$. The positive sign follows from the two negative cross terms.
The Laplace identity follows by the spectral theorem and
Theorem~\ref{thm:schur}. This explicitly realizes projection memory of
the kind introduced in \cite{Mori}, with this Euclidean generator.

For every smooth $f$, orthogonality of the two components of
$\LL J_mf$ gives the full second moment
\begin{equation}\label{eq:second}
 \|\LL J_mf\|^2=\|A_mf\|^2+\|B_mf\|^2.
\end{equation}
Thus $K_m(t)$ is replaced by $e^{-tA_m}$ for all $t$ only when $B_m$
vanishes: compare their scalar second derivatives at $t=0$ on a smooth
core. Conversely $B_m=0$ makes \eqref{eq:block} a direct-sum form and
proves that replacement exactly. We do not make that replacement.

The dependence on energy also records the actual reconstructed norm.
For fixed $0\ne f\in\VV_m$ with $\mu_m(f)=0$ and scalar $z\geq0$, put
\begin{equation}\label{eq:dressed}
 w_z=(C_m+z)^{-1}B_mf,\quad
 u_z=\psi(J_mf-w_z),\quad d_f=\|f\|^2.
\end{equation}
The $z=0$ case exists because $C_m\geq\delta>0$ in every fixed box.
The vector is in the original form domain, is orthogonal to $\psi$,
and is gauge invariant when $f$ is. Its exact norm and energy are
\begin{align}
 \|u_z\|_\lambda^2&=d_f+\|w_z\|_\mu^2
                  =\partial_z s_{m,z}(f,f),\label{eq:normderivative}\\
 \mathfrak q_{H-E_0}[u_z]
 &=a_m(f,f)-\langle B_mf,w_z\rangle-z\|w_z\|^2
   =s_{m,z}(f,f)-z\partial_zs_{m,z}(f,f).\label{eq:dressedenergy}
\end{align}
To verify these, orthogonality gives the norm; differentiation of the
bounded resolvent gives $\partial_z(C_m+z)^{-1}=-(C_m+z)^{-2}$.
Testing $(C_m+z)w_z=B_mf$ against $w_z$ gives
$c_m(w_z,w_z)=\langle B_mf,w_z\rangle-z\|w_z\|^2$; insert it in
\eqref{eq:block}. Nonnegativity of $q$ proves the numerator is nonnegative.
For $B_mf\ne0$ the quotient in \eqref{eq:dressedenergy} divided by
\eqref{eq:normderivative} is strictly below $a_m(f,f)/d_f$.
These are energies of actual full-system vectors; no artificial effective
vacuum norm has replaced their eliminated components.

\section{Associativity of repeated blocking with all prior memory}

Let $m\mid n\mid2L$. A coarse $n$-edge is the ordered concatenation of
$n/m$ coarse $m$-edges along its line. Let $\pi_{m,n}$ take these products
and discard the other $m$-coarse edges. Literal ordered multiplication gives
\begin{equation}\label{eq:tower}
 \pi_n=\pi_{m,n}\circ\pi_m,\quad J_n=J_m I_{m,n},\quad
 \mu_n=(\pi_{m,n})_*\mu_m,\quad
 J_n^*=I_{m,n}^*J_m^*.
\end{equation}
Here $I_{m,n}f=f\circ\pi_{m,n}$ is an isometry $\KK_n\to\KK_m$;
its adjoint is conditional integration using $r_m$, not Haar integration
in place of that weight. Fubini proves each equality, including the adjoint
identity. In the full Hilbert space $P_nP_m=P_mP_n=P_n$.
All maps preserve the relevant Sobolev spaces by the same compact
path-product coordinate construction. Coefficient composition in the
instantaneous forms is exactly $\kappa m(n/m)=\kappa n$.

Let $\widehat Q=I-I_{m,n}I_{m,n}^*$ on $\KK_m$. On
$\widehat Q\KK_m\cap\VV_m$, let $t_{22,z}$ be the restriction of the
already completed form $s_{m,z}$. Define its other block forms by
\begin{align*}
 t_{11,z}(f,g)&=s_{m,z}(I_{m,n}f,I_{m,n}g),\\
 t_{21,z}(h,f)&=s_{m,z}(h,I_{m,n}f),\\
 t_{12,z}(f,h)&=s_{m,z}(I_{m,n}f,h).
\end{align*}
The operator $T_{22,z}$ associated to $t_{22,z}$ is positive with lower
bound $z$. Its inverse also acts from its form dual to its form space,
by the coercive weak equation. In that sense the next eliminated component
is $h_f=-T_{22,z}^{-1}t_{21,z}(\cdot,f)$, and the next form is
\begin{equation}\label{eq:recursive}
 s_{n,z}(f,g)=t_{11,z}(f,g)
       -t_{12,z}\bigl(f,T_{22,z}^{-1}t_{21,z}(\cdot,g)\bigr).
\end{equation}
This notation means substitution of the unique weak solution, so does
not require multiplying unspecified unbounded compressions.

\begin{theorem}\label{thm:associative}
Equation \eqref{eq:recursive} equals the direct fine-to-$n$ form
\eqref{eq:schur}. In particular
\begin{equation}\label{eq:resolventtower}
 S_{n,z}^{-1}=I_{m,n}^*S_{m,z}^{-1}I_{m,n}
            =J_n^*(\LL+z)^{-1}J_n.
\end{equation}
No memory generated at an earlier level is lost in the next level.
\end{theorem}
\begin{proof}
For prescribed $f\in\VV_n$, every $u\in\VV$ with $J_n^*u=f$ has a
unique orthogonal decomposition
$u=J_n f+J_mh+k$, where $h\in\widehat Q\KK_m\cap\VV_m$ and
$k\in\HH_{Q_m}\cap\VV$; obtain it by the nested projections in
\eqref{eq:tower}. Put $v_h=J_m(I_{m,n}f+h)$ to keep both arguments
of the sesquilinear form explicit. The exact variational identity
\eqref{eq:square} gives
\begin{align*}
 s_{n,z}(f,f)
 &=\inf_{J_n^*u=f}[q(u,u)+z\|u\|^2]\\
 &=\inf_h\inf_k[q(v_h+k,v_h+k)+z\|v_h+k\|^2]\\
 &=\inf_h s_{m,z}(I_{m,n}f+h,I_{m,n}f+h).
\end{align*}
The infima have unique minimizers by coercivity and the closed form-space
constraints. The weak equation for the last minimizer is precisely the
$T_{22,z}$ equation above. Expanding the last form and polarizing proves
\eqref{eq:recursive}. Alternatively compress
\eqref{eq:resolvent} by $I_{m,n}$ and use \eqref{eq:tower}; this proves
\eqref{eq:resolventtower} directly and confirms the same unique closed
form. Both arguments include all intermediate eliminated components.
\end{proof}

To display the generated terms explicitly, write
$\mathcal M_{m,z}(f,g)=\langle B_mf,(C_m+z)^{-1}B_mg\rangle$.
Then each of $t_{11,z},t_{12,z},t_{21,z},t_{22,z}$ is the corresponding
block of $a_m+z\langle\cdot,\cdot\rangle-\mathcal M_{m,z}$.
In particular the two off-diagonal blocks contain the corresponding
off-diagonal terms of $-\mathcal M_{m,z}$, and $T_{22,z}$ contains its
entire diagonal term. Setting these terms to zero would change
\eqref{eq:recursive}. The exact recursion is independent of the order of
the specified nested eliminations; it does not reinitialize the dynamics
at each marginal $A_m$.

\section{A large-block Wilson channel with uniform norm and spectral bounds}

Let $C$ be a closed coarse edge word using each coarse physical edge at
most once. Its fine path then uses each of its fine physical edges once;
let $r$ be its number of coarse edges and $\ell=mr$ its number of fine
edges. For example, every elementary coarse face has $r=4$, fine perimeter
$\ell=4m$, and physical side length $ma$. Keep its original ordered word,
with inverses on the negative traversals, and set
\begin{gather}
 W_C(z)=\tr\prod_{(\gamma,\epsilon)\in C}z_\gamma^\epsilon,\quad
 \overline W_C=\mu_m(W_C),\quad r_C=\mu_m(W_C^2),\nonumber\\
 f_C=W_C-\overline W_C,\quad u_C=\psi J_mf_C,\quad
 d_C=\|u_C\|_\lambda^2=r_C-(\overline W_C)^2.\label{eq:loop}
\end{gather}
SU(2) traces are real. The fine and coarse gauge actions both fix the
closed trace, so $u_C$ is a physical invariant vector, orthogonal to
$\psi$. The subtraction is the actual interacting mean.

\begin{lemma}\label{lem:loopmoments}
For every such loop, at arbitrary positive coupling,
\begin{equation}\label{eq:loopnorm}
 \alpha\leq d_C\leq4,\qquad
 \alpha\leq r_C\leq4-3\alpha.
\end{equation}
The exact unnormalized excitation energy and bounds are
\begin{equation}\label{eq:loopenergy}
 \mathfrak q_{H-E_0}[u_C]
   =\kappa\ell\left(1-\frac{r_C}{4}\right),\qquad
 \frac34\alpha\kappa\ell\leq\mathfrak q_{H-E_0}[u_C]\leq\kappa\ell.
\end{equation}
Consequently
\begin{equation}\label{eq:loopquotient}
 \frac{3\alpha}{16}\kappa\ell\leq
 \frac{\mathfrak q_{H-E_0}[u_C]}{d_C}
 =\kappa\ell\frac{1-r_C/4}{r_C-(\overline W_C)^2}
 \leq\frac{\kappa\ell}{\alpha}.
\end{equation}
\end{lemma}
\begin{proof}
Fix any one fine edge used by the loop and condition on all other fine
variables. The ordered holonomy is $A u^{\pm1}B$ for fixed group matrices
$A,B$, so its Haar distribution is Haar by invariance and inversion.
Writing $U=u_0I+i\sum_a u_a\sigma_a$ identifies Haar measure with the
uniform measure on the unit sphere in $\mathbb R^4$. Its coordinate
second moments are $1/4$ by sign and permutation symmetry and
$\sum_{a=0}^3u_a^2=1$. Thus
$\int W_C\dd u=0$, $\int W_C^2\dd u=1$, and
$\int(1-W_C^2/4)\dd u=3/4$.
The conditional density is bounded below by $\alpha$ from
Theorem~\ref{thm:localvacuum}. Integrating the two nonnegative functions
$W_C^2$ and $4-W_C^2$ proves the two bounds on $r_C$.
For every real constant $c$, conditional integration gives
$\int(W_C-c)^2p(u\mid y)\dd u\geq\alpha\int(W_C-c)^2\dd u
=\alpha(1+c^2)\geq\alpha$.
Taking $c$ to be that conditional mean proves conditional variance
at least $\alpha$. The law of total variance, obtained by expanding
around the conditional mean and integrating the cross term to zero,
proves $d_C\geq\alpha$. Since $|W_C|\leq2$, $d_C\leq4$.

For a positive occurrence, a based form of the derivative is
$\tr(T_aA_C)$ with $A_C\in\SU(2)$ having trace $W_C$; for a negative
occurrence it is $-\tr(T_aA_C')$. This follows from
$X_a u^{-1}=-u^{-1}T_a$ and cyclic trace invariance. The Pauli expression
above gives $\sum_a|\tr(T_aA_C)|^2=1-W_C^2/4$.
There are exactly $\ell$ used fine edges. Insert this identity in
\eqref{eq:q}; centering leaves derivatives unchanged. The exact energy
and its lower and upper bounds follow. Divide using
$\alpha\leq d_C\leq4$ to obtain \eqref{eq:loopquotient}.
\end{proof}

These bounds are for the same density used by the block projection.
For an additional consistency check, that coarse density itself has a
single-coarse-link conditional lower bound $\alpha$, independent of $m$.
Indeed change only the last fine edge of a specified chain, holding all
$y$ and all other coarse links fixed in \eqref{eq:haarcoordinates}.
Theorem~\ref{thm:localvacuum} bounds the ratio of the two full densities
by $R^2$. Integrating the same $y$ on both sides bounds the ratio
$r_m(z_\gamma,z_{\ne\gamma})/r_m(z'_\gamma,z_{\ne\gamma})$ by $R^2$.
Averaging over $z'_\gamma$ proves the assertion. This exact argument
does not equate single-link conditional bounds with independence or a
global Poincare inequality.

\begin{theorem}\label{thm:spectral}
Define the unnormalized actual full-Hamiltonian spectral measure
\begin{equation}\label{eq:nu}
 \mathsf M_C(B)=\langle u_C,\mathbf1_B(H-E_0)u_C\rangle,\qquad
 \mathsf M_C([0,\infty))=d_C.
\end{equation}
No division by $d_C$ is made in this measure. It has no atom at zero.
Its full first and second moments satisfy
\begin{gather}
 \int\omega\dd\mathsf M_C=\kappa\ell(1-r_C/4),\nonumber\\
 \int\omega^2\dd\mathsf M_C
 =\|A_mf_C\|^2+\|B_mf_C\|^2
 \leq C_*^2\kappa^2\ell^2.\label{eq:moments}\\
 \Omega_C=\frac{3\alpha}{32}\kappa\ell,\qquad
 c_H=\frac{9\alpha^3}{1024C_*^2}>0,\qquad
 \boxed{\mathsf M_C([\Omega_C,\infty))\geq c_Hd_C.}\label{eq:highweight}
\end{gather}
In particular $c_H$ is independent of the increasing blocking length $m$
and of total box volume. The full compressed resolvent, including its
memory, obeys at the scalar shift $z=\kappa\ell$
\begin{equation}\label{eq:resolventdeficit}
 z\langle f_C,S_{m,z}^{-1}f_C\rangle
 \leq d_C\left[1-c_H\frac{3\alpha/32}{1+3\alpha/32}\right].
\end{equation}
\end{theorem}
\begin{proof}
The state is smooth and orthogonal to the unique vacuum. Spectral
calculus and \eqref{eq:loopenergy} prove the first moment and the
absence of a zero atom. Every used edge has fundamental Casimir $3/4$,
including inverse occurrences, so
$-\sum_{e,a}X_{e,a}^2W_C=3\ell W_C/4$.
Combining \eqref{eq:gs} with \eqref{eq:scorebound} and
$|\nabla_e W_C|\leq1$ gives the pointwise bound
\[
 |\LL J_mf_C|=|\LL J_mW_C|
 \leq\kappa\ell\left(\frac32+2G\right)=\kappa\ell C_*.
\]
The second-moment bound follows by integration and
\eqref{eq:second}. In particular the term $\|B_mf_C\|^2$ has not been
deleted from this spectral measure.

Put $A=3\alpha/16$. The first moment is at least $A\kappa\ell d_C$
by \eqref{eq:loopquotient}. With $\Omega_C=A\kappa\ell/2$, its part on
$[\Omega_C,\infty)$ is at least $A\kappa\ell d_C/2$, since the integral
over $[0,\Omega_C)$ is at most $\Omega_C d_C$.
Cauchy--Schwarz against the unnormalized $\mathsf M_C$ gives
\[
 (A\kappa\ell d_C/2)^2
 \leq\left(\int\omega^2\dd\mathsf M_C\right)
                       \mathsf M_C([\Omega_C,\infty)).
\]
Using \eqref{eq:moments} gives the stronger intermediate bound
$\mathsf M_C([\Omega_C,\infty))\geq A^2d_C^2/(4C_*^2)$.
Since $d_C\geq\alpha$, this is at least
$A^2\alpha d_C/(4C_*^2)=c_Hd_C$, with the state mass retained.
Finally \eqref{eq:resolvent} gives
$z\langle f_C,S_{m,z}^{-1}f_C\rangle
=\int z/(\omega+z)\dd\mathsf M_C$.
On $[\Omega_C,\infty)$ its integrand is at most
$z/(\Omega_C+z)$, and elsewhere it is at most $1$.
Its upper bound is therefore
$d_C-\Omega_C\mathsf M_C([\Omega_C,\infty))/(\Omega_C+z)$.
Insert \eqref{eq:highweight} and $z=\kappa\ell$ to obtain
\eqref{eq:resolventdeficit}.
\end{proof}

The same channel can be dressed by the complete eliminated resolvent
with a quantitative bound that continues to hold for growing $m$.
Define
\begin{equation}\label{eq:tau}
 \tau_* =\frac{8G}{\sqrt{\alpha c_H}},\qquad
 z_C=\tau\kappa\ell,\quad\tau\geq\tau_*.
\end{equation}
For $f=f_C$ use the full $u_{z_C}$ of \eqref{eq:dressed}, with all
fine eliminated degrees of freedom in $C_m$.
\begin{corollary}\label{cor:dressedweight}
This is a nonzero physical vector with exact norm and energy
\eqref{eq:normderivative}--\eqref{eq:dressedenergy}. Quantitatively,
\begin{gather}
 d_C\leq\|u_{z_C}\|^2\leq d_C+\alpha c_H/4,\nonumber\\
 \|\mathbf1_{[\Omega_C,\infty)}(H-E_0)u_{z_C}\|^2
 \geq\frac{c_H}{4+c_H}\|u_{z_C}\|^2,\label{eq:dressedweight}\\
 \frac{\Omega_Cc_H}{4+c_H}
 \leq\frac{\mathfrak q_{H-E_0}[u_{z_C}]}{\|u_{z_C}\|^2}
 \leq\kappa\ell\frac{1-r_C/4}{d_C}.
\end{gather}
\end{corollary}
\begin{proof}
By \eqref{eq:coarsescorebound}, each used coarse edge contributes at
most $4G\kappa m$ to the pointwise absolute value of $B_mf_C$;
there are $r$ such edges and $|D_\gamma W_C|\leq1$.
Thus $\|B_mf_C\|\leq4G\kappa\ell$.
Since $C_m\geq0$, its shifted inverse has norm at most $1/z_C$.
Consequently
\[
 \|w_{z_C}\|\leq4G/\tau\leq\tfrac12\sqrt{\alpha c_H}
                              \leq\tfrac12\sqrt{d_Cc_H}.
\]
Orthogonality gives the norm bound. Multiplication by $\psi$ is unitary,
and an orthogonal spectral projection is contractive. Hence the triangle
inequality and \eqref{eq:highweight} give
\[
 \|\mathbf1_{[\Omega_C,\infty)}(H-E_0)u_{z_C}\|
 \geq\sqrt{d_Cc_H}-\|w_{z_C}\|
 \geq\tfrac12\sqrt{d_Cc_H}.
\]
The retained norm satisfies
$\|u_{z_C}\|^2=d_C+\|w_{z_C}\|^2\leq d_C(1+c_H/4)$,
so the squared lower bound $d_Cc_H/4$ is at least
$c_H\|u_{z_C}\|^2/(4+c_H)$.
Integrating $\omega\geq\Omega_C$ over this weight proves the lower
energy bound; \eqref{eq:dressedenergy} proves the upper bound.
The large positive shift here is explicitly fixed in \eqref{eq:tau};
this argument does not assert the same lower bound for the separately
defined vector minimizing the energy numerator at $z=0$. Its full
Rayleigh variational problem is calculated in Section~\ref{sec:zero}.
\end{proof}


\section{Zero-shift reconstruction and affine physical states}
\label{sec:zero}

We now calculate the zero-shift channel without transferring the
positive-shift estimate by assertion. Gauge invariance is imposed
throughout this section. Let $\HH^{\rm G}$ be the closed subspace of
gauge-invariant functions in $\HH$, and let $\KK_m^{\rm G}$ be the
corresponding coarse subspace. The gauge group is the compact product
$\SU(2)^{\mathsf V_L}$; its Haar average is an orthogonal projection.
The intertwining of the path products and the gauge descent of the
conditional density show that $J_m,J_m^*,P_m,Q_m$ commute with the
corresponding gauge actions. Consequently $B_m,C_m$ and their
resolvents preserve these subspaces. All physical vectors below are
obtained by the same unitary multiplication by $\psi$.

Take $m\geq2$ dividing $2L$. In this section write $C$ for the restriction
of $C_m$ to $\HH_{Q_m}^{\rm G}$, retaining precisely its form domain
$\VV_{Q_m}\cap\HH^{\rm G}$. This is a reducing restriction of the
operator, not a change of interaction. Put
\[
 c=\inf\sigma(C),\qquad
 \delta_{\rm G}=\inf_{\substack{0\ne u\in\VV\cap\HH^{\rm G}\\\mu u=0}}
                   \frac{q(u,u)}{\|u\|^2}.
\]
This complementary physical space is nonzero. For example, the fine
face at $n=(-L,-L,-L)$ in the first two directions contains the
direction-2 edge with first coordinate $-L+1$, which is outside $R_m$.
Varying that edge with every other fine edge fixed leaves every coarse
product unchanged and changes the fine face trace. Positivity of the
conditional density implies $Q_mW_p\ne0$, and this function is gauge
invariant. Compact embedding of the restricted form domain into $L^2$
gives compact resolvent for $C$: a bounded sequence in its form norm
is bounded in the original $H^1$ norm. Thus $c$ is an eigenvalue.

\begin{proposition}\label{prop:finitezero}
With the original number of fine edges $N=3(2L)(2L+1)^2$, define
\begin{equation}\label{eq:finitegap}
 \Delta_L=\frac{3\kappa}{4}\alpha^N
          =\frac{3g^2}{2a}\alpha^{\,3(2L)(2L+1)^2}.
\end{equation}
Then $\delta_{\rm G}\geq\delta\geq\Delta_L>0$ and $c\geq\delta_{\rm G}$.
For mean-zero $f\in\VV_m\cap\KK_m^{\rm G}$, define the linear map
\begin{equation}\label{eq:Rzero}
 \mathcal R_m f=J_mf-C^{-1}B_mf,\quad
 s_{m,0}(f,g)=a_m(f,g)-\langle B_mf,C^{-1}B_mg\rangle.
\end{equation}
The form $s_{m,0}$ is closed on the indicated retained form space,
with
\begin{equation}\label{eq:zerocoercivity}
 \Delta_L\|f\|^2\leq s_{m,0}(f,f)\leq a_m(f,f),\qquad
 a_m(f,f)\leq2(1+\eta_m/\Delta_L)s_{m,0}(f,f).
\end{equation}
The map $\mathcal R_m$ takes this space into the original physical
mean-zero form domain and satisfies $J_m^*\mathcal R_mf=f$.
For every $h\in\VV_{Q_m}\cap\HH^{\rm G}$,
\begin{equation}\label{eq:zeroorthogonal}
 q(\mathcal R_m f,h)=0,\qquad
 q(\mathcal R_m f,\mathcal R_m g)=s_{m,0}(f,g).
\end{equation}
For its associated mean-zero operator $S_{m,0}$, the exact domain map is
\begin{equation}\label{eq:zerodomain}
 f\in D(S_{m,0})\ \Longleftrightarrow
 \mathcal R_m f\in D(\LL),\qquad
 \LL\mathcal R_m f=J_mS_{m,0}f.
\end{equation}
The equivalence in \eqref{eq:zerodomain} is for $f$ already in the
retained mean-zero form domain.
\end{proposition}
\begin{proof}
Change fine variables one at a time between any two configurations.
Theorem~\ref{thm:localvacuum} gives
$\rho_{\max}/\rho_{\min}\leq R^{2N}=\alpha^{-N}$.
The product Haar Laplacian has first positive eigenvalue $3/4$:
tensor products of the complete matrix-coefficient bases in
Lemma~\ref{lem:singleheat} have eigenvalues
$\sum_e j_e(j_e+1)$, whose least positive value is $3/4$.
Expanding a smooth function in that basis, and then using $H^1$
density, gives the Haar Poincare inequality. Hence
\[
 \Var_\mu u
 \leq\rho_{\max}\Var_\lambda u
 \leq\frac{4\rho_{\max}}3\sum_{e,a}\int|X_{e,a}u|^2\dd\lambda
 \leq\frac{4\rho_{\max}}{3\kappa\rho_{\min}}q(u,u).
\]
The first inequality evaluates the infimum over constants defining
$\Var_\mu$ at the Haar mean. This proves \eqref{eq:finitegap} for all
mean-zero functions, and restricting to gauge-invariant functions
preserves the lower bound. Every complementary physical function has
zero mean, so $c\geq\delta_{\rm G}$.

The bounded resolvent $C^{-1}$ maps its Hilbert space to $D(C)$ and
therefore to the form domain. Completing the square at zero gives
\[
 q(J_mf+h,J_mf+h)
 =s_{m,0}(f,f)+c_m(h+C^{-1}B_mf,h+C^{-1}B_mf).
\]
The unique minimizer is $\mathcal R_m f$, proving
\eqref{eq:zeroorthogonal} and the upper bound in
\eqref{eq:zerocoercivity}. Its mean is zero and its squared norm is
$\|f\|^2+\|C^{-1}B_mf\|^2$. The already proved full Poincare inequality
gives the lower bound and
$\|\mathcal R_m f\|^2\leq s_{m,0}(f,f)/\Delta_L$.
Apply \eqref{eq:projectionbound} to $\mathcal R_m f$ for the second
inequality in \eqref{eq:zerocoercivity}. These two-sided form-norm
bounds prove closedness with exactly the retained $H^1$ domain.

If $f\in D(S_{m,0})$, decompose any physical form test vector as
$J_mg+h$. Equations \eqref{eq:zeroorthogonal} give
\[
 q(J_mg+h,\mathcal R_m f)=s_{m,0}(g,f)
                       =\langle J_mg+h,J_mS_{m,0}f\rangle .
\]
Constants give zero on both sides. This is the defining weak operator
identity. It also holds against arbitrary, noninvariant test vectors:
average that test vector over the gauge group, since the operator and
the vector being tested are invariant. Thus $\mathcal R_m f$ belongs
to the full $D(\LL)=H^2$ and has the displayed image.
Conversely if $\mathcal R_m f\in D(\LL)$, tests in the complementary
space show $Q_m\LL\mathcal R_m f=0$; retained tests then prove
$f\in D(S_{m,0})$ and the claimed image.
\end{proof}

The exponential dependence on $N$ in \eqref{eq:finitegap} is explicit.
It provides a proved finite-box inverse bound, not a volume-uniform gap.

Fix a nonzero, mean-zero, gauge-invariant $f\in\VV_m$. Keep its norm
unchanged and set
\[
 d_f=\|f\|^2,\quad a_f=a_m(f,f),\quad b_f=B_mf,\quad
 \sigma_f(I)=\|\mathbf1_I(C)b_f\|^2 .
\]
This $\sigma_f$ is a finite positive measure on $[c,\infty)$ of total
mass $\|b_f\|^2$, with no division by that mass. For every real $z>-c$,
put
\begin{equation}\label{eq:inversemoments}
 M_j(z)=\int_{[c,\infty)}(t+z)^{-j}\dd\sigma_f(t),\quad j=1,2,3,
 \qquad v_z=J_mf-(C+z)^{-1}b_f .
\end{equation}
All these moments are finite by $t+z\geq c+z>0$.

\begin{theorem}\label{thm:affine}
The actual physical state $\psi v_z$ has exact squared norm, excitation
energy and Rayleigh quotient
\begin{equation}\label{eq:affinemoments}
 N_f(z)=d_f+M_2(z),\quad
 e_f(z)=a_f-M_1(z)-zM_2(z),\quad
 Q_f(z)=\frac{e_f(z)}{N_f(z)}.
\end{equation}
For all $z>-c$,
\begin{equation}\label{eq:rayleighderivative}
 Q_f'(z)=\frac{2M_3(z)}{N_f(z)}[z+Q_f(z)].
\end{equation}
In particular $Q_f$ is strictly increasing on $[0,\infty)$ whenever
$b_f\ne0$. Its zero-shift value is strictly smaller than the lifted
quotient $a_f/d_f$, and some negative shifts give a quotient strictly
smaller still. Zero shift minimizes the energy numerator at fixed
retained data, but it does not minimize the physical quotient when
$b_f\ne0$.

The full affine variational problem, with all physical eliminated
states allowed, is determined as follows:
\begin{gather}
 \chi_f=\inf_{h\in\VV_{Q_m}\cap\HH^{\rm G}}
           \frac{q(J_mf+h,J_mf+h)}{d_f+\|h\|^2},\nonumber\\
 F_f(E)=a_f-E d_f-M_1(-E),\qquad 0\leq E<c,\qquad
 F_f'(E)=-d_f-M_2(-E)<0,\label{eq:affineroot}\\
 F_f(0)>0,\quad
 \mathfrak l_f=\lim_{E\uparrow c}F_f(E)\in[-\infty,\infty).\nonumber
\end{gather}
If $\mathfrak l_f<0$, there is exactly one root $E_f\in(0,c)$ and
\begin{gather}
 \chi_f=E_f,\qquad
 h_f=-(C-E_f)^{-1}b_f,\qquad
 \|\psi(J_mf+h_f)\|^2=-F_f'(E_f),\nonumber\\
 \mathfrak q_{H-E_0}[\psi(J_mf+h_f)]=-E_fF_f'(E_f).
 \label{eq:affineattain}
\end{gather}
This is the unique affine minimizer. If $\mathfrak l_f\geq0$, then
$\chi_f=c$. In this case $b_f$ is orthogonal to $\ker(C-c)$.
The value $c$ is attained exactly when $\mathfrak l_f=0$, and all
minimizers are
\begin{equation}\label{eq:boundaryattain}
 h=-(C-c)^\dagger b_f+k,\qquad k\in\ker(C-c).
\end{equation}
Here the dagger is the inverse on $\ker(C-c)^\perp$ and zero on the
kernel. When $\mathfrak l_f>0$, no finite affine vector attains $c$;
the infimum is approached by $h=t k$, $t\to\infty$, for a nonzero
$k\in\ker(C-c)$. These alternatives are exhaustive for the actual
interacting $C$ and $b_f$, including $b_f=0$.
\end{theorem}
\begin{proof}
The spectral theorem identifies $M_1$ with
$\langle b_f,(C+z)^{-1}b_f\rangle$ and $M_2$ with
$\|(C+z)^{-1}b_f\|^2$. The orthogonal decomposition of $v_z$ and
the block form \eqref{eq:block} prove \eqref{eq:affinemoments}
for negative as well as positive $z$. Differentiating these bounded
resolvents gives $M_1'=-M_2$, $M_2'=-2M_3$,
$N_f'=-2M_3$, and $e_f'=2zM_3$. The quotient rule proves
\eqref{eq:rayleighderivative}. Each $v_z$ is nonzero and mean-zero,
so $Q_f(z)\geq\delta_{\rm G}>0$.
For $b_f\ne0$, $M_3>0$ and the derivative is positive at zero
and at every positive shift. Also
\[
 \frac{a_f}{d_f}-Q_f(0)
 =\frac{M_1(0)+(a_f/d_f)M_2(0)}{d_f+M_2(0)}>0 .
\]
Continuity of the derivative at zero proves strict improvement for
sufficiently small negative shifts. The statement about the numerator
follows independently from \eqref{eq:zeroorthogonal}; increasing the
norm is part of the exact quotient calculation.

For $0\leq E<c$, complete the shifted square for arbitrary $h$:
\begin{equation}\label{eq:affinesquare}
 q(J_mf+h,J_mf+h)-E(d_f+\|h\|^2)
 =F_f(E)+\|(C-E)^{1/2}[h+(C-E)^{-1}b_f]\|^2.
\end{equation}
No discarded term or approximation enters this identity.
The sign and derivative statements in \eqref{eq:affineroot} follow
from Proposition~\ref{prop:finitezero} and the spectral derivative.
Monotonicity bounds the limit above by $F_f(0)$ and allows the value
$-\infty$ below. The intermediate value theorem proves the unique
root when the limit is negative. At this root
\eqref{eq:affinesquare} bounds every quotient below by $E_f$,
and equality holds precisely for the stated $h_f$. The norm and
energy in \eqref{eq:affineattain} follow from
\eqref{eq:affinemoments} and $F_f(E_f)=0$.

If the limit is nonnegative, $F_f(E)>0$ for every $E<c$.
Equation \eqref{eq:affinesquare}, followed by $E\uparrow c$,
bounds every affine quotient below by $c$. A complementary
eigenvector $k$ of eigenvalue $c$ yields
\[
 \frac{a_f+2t\Rea\langle b_f,k\rangle+t^2c\|k\|^2}
      {d_f+t^2\|k\|^2}\longrightarrow c,
\]
proving equality of the infimum. If $b_f$ had a nonzero projection
on that eigenspace, its contribution to $M_1(-E)$ would diverge as
$(c-E)^{-1}$, contradicting $\mathfrak l_f\geq0$.
Compact resolvent isolates $c$ from the rest of $\sigma(C)$, so
$(C-c)^\dagger$ is bounded on the complementary eigenspaces and
maps into $D(C)$. Completing the square at $E=c$ now gives
\[
 q(J_mf+h,J_mf+h)-c(d_f+\|h\|^2)
 =\mathfrak l_f+
 \|(C-c)^{1/2}[h+(C-c)^\dagger b_f]\|^2 .
\]
It proves both the attainment criterion and every vector in
\eqref{eq:boundaryattain}.
\end{proof}

The scalar root in \eqref{eq:affineroot} solves the affine problem for
this fixed retained direction. It is not, in general, an eigenvalue of
the full operator: a full eigenvector also requires the weak equation
$s_{m,-E}(g,f)=0$ for every retained test direction $g$.
Here $s_{m,-E}$ denotes the sesquilinear expression
\eqref{eq:schur} at $z=-E$, restricted to physical retained functions;
it is well defined on $\VV_m\cap\KK_m^{\rm G}$ because $E<c$.
Indeed complementary testing gives $h=-(C-E)^{-1}B_mf$, and retained
testing gives exactly that additional equation. Conversely these
equations imply the full weak eigenvector equation and its $H^2$
operator domain. This proves the exact map in both directions, and
explains precisely which test directions the affine minimization uses.
The energy-dependent eigenvalue and reconstruction principle agrees
with \cite[Theorem~1.2 and Corollary~1.3]{DSS}; the affine classification
and the physical domains above were proved directly.

\begin{corollary}\label{cor:zeroloop}
For the loop $f=f_C$, the zero-shift state
$u_0=\psi\mathcal R_m f_C$ obeys the exact identities
\begin{align}
 \|u_0\|^2&=d_C+M_2(0),\nonumber\\
 \mathfrak q_{H-E_0}[u_0]&=\kappa\ell(1-r_C/4)-M_1(0).
 \label{eq:zeroloopactual}
\end{align}
\begin{gather}
 \alpha\leq d_C\leq\|u_0\|^2
 \leq d_C+\min\left\{
       \frac{\kappa\ell(1-r_C/4)}{\Delta_L},
       \frac{16G^2\kappa^2\ell^2}{\Delta_L^2}\right\},
 \label{eq:zeroloopnorm}\\
 \|u_0\|^2\leq4+\frac{4\ell}{3\alpha^N},\nonumber\\
 \Delta_L\leq
 \frac{\mathfrak q_{H-E_0}[u_0]}{\|u_0\|^2}
 \leq\frac{\kappa\ell(1-r_C/4)}{d_C}.\label{eq:zeroloopenergy}
\end{gather}
For an elementary coarse face, the last norm bound is
$4+16m/(3\alpha^N)$, and its energy numerator is exactly
$8g^2m(1-r_C/4)/a-M_1(0)$. The same $\Delta_L=3g^2\alpha^N/(2a)$
bounds the quotient below. These bounds are finite and explicit;
their volume dependence is retained.
\end{corollary}
\begin{proof}
The identities are \eqref{eq:affinemoments} at zero.
Positivity gives $M_1(0)\leq a_f$, and the spectral inequality
$t^{-2}\leq c^{-1}t^{-1}$ gives
$M_2(0)\leq a_f/c\leq a_f/\Delta_L$.
The other estimate uses $M_2(0)\leq\|B_mf_C\|^2/\Delta_L^2$
and the already proved $\|B_mf_C\|\leq4G\kappa\ell$.
Insert $a_f=\kappa\ell(1-r_C/4)\leq\kappa\ell$ and
$d_C\leq4$. The lower quotient bound is the full Poincare inequality,
and the upper one follows from the exact numerator and norm.
Finally use $\ell=4m$ and $\kappa=2g^2/a$.
\end{proof}

\section{Exact comparison of all physical low-energy states}
\label{sec:allsoft}

The previous section treated a fixed retained observable. We now give
a quantitative relation encompassing all physical states, with no
assumed complementary gap uniform in volume. Define the actual
zero-shift reconstructed infimum
\begin{equation}\label{eq:reconstructedgap}
 \lambda_m^{\rm R}
 =\inf_{\substack{0\ne f\in\VV_m\cap\KK_m^{\rm G}\\\mu_m f=0}}
 \frac{s_{m,0}(f,f)}
      {\|f\|^2+\|C^{-1}B_mf\|^2}.
\end{equation}
This denominator is the full physical norm. In every fixed box it is
positive, and $\lambda_m^{\rm R}\geq\delta_{\rm G}\geq\Delta_L$.

\begin{theorem}\label{thm:gaptransfer}
For every allowed $L,m,a,g$ with $m\geq2$, the following comparison
holds with numerical constants independent of all four parameters:
\begin{equation}\label{eq:harmonicgap}
 \frac{\lambda_m^{\rm R}c}{\lambda_m^{\rm R}+c}
 \leq\delta_{\rm G}\leq\min\{\lambda_m^{\rm R},c\},\qquad
 \frac12\min\{\lambda_m^{\rm R},c\}\leq\delta_{\rm G}.
\end{equation}
It has the following explicit state map. For every mean-zero physical
$u\in\VV$, put
\begin{equation}\label{eq:softdecomposition}
 f=J_m^*u,\quad v=\mathcal R_m f,\quad
 k=u-v=Q_mu+C^{-1}B_mf .
\end{equation}
Then $k\in\VV_{Q_m}\cap\HH^{\rm G}$ and
\begin{equation}\label{eq:softidentities}
 q(u,u)=s_{m,0}(f,f)+c_m(k,k),\qquad
 \|u\|^2=\|v\|^2+\|k\|^2+2\Rea\langle v,k\rangle .
\end{equation}
At least one of the nonzero components $v,k$ has physical Rayleigh
quotient at most $2q(u,u)/\|u\|^2$. Conversely each such component is an
actual admissible physical state under multiplication by $\psi$.

For every sequence of allowed boxes, block lengths, positive spacings
and couplings, $\delta_{\rm G}$ tends to zero if and only if
$\min\{\lambda_m^{\rm R},c\}$ tends to zero. Uniform positive lower
bounds for $\delta_{\rm G}$ are equivalent to simultaneous uniform
positive lower bounds for these two actual sector infima. These are
equivalences proved by \eqref{eq:harmonicgap}, not assumed estimates
for either sector.
\end{theorem}
\begin{proof}
Sobolev preservation of $J_m^*$ and Proposition~\ref{prop:finitezero}
make all components in \eqref{eq:softdecomposition} well defined in
their stated form domains. The identity $J_m^*v=f$ puts $k$ in the
complement. Equation \eqref{eq:zeroorthogonal} gives $q(v,k)=0$,
which proves the energy identity; the norm identity retains its
cross term. Cauchy--Schwarz in two scalar components gives
\[
 \|u\|\leq\|v\|+\|k\|
 \leq\sqrt{\frac{q(v,v)}{\lambda_m^{\rm R}}}
      +\sqrt{\frac{c_m(k,k)}c}
 \leq\sqrt{\left(\frac1{\lambda_m^{\rm R}}+\frac1c\right)q(u,u)}.
\]
Zero components contribute zero. Taking the infimum over $u$ proves
the first lower bound. The upper bound follows because the full
variational set contains both the reconstructed vectors and all
complementary physical vectors. For positive $x,y$,
$xy/(x+y)\geq\min(x,y)/2$, proving the last bound.

For the explicit factor-two state assertion,
$\|u\|^2\leq2(\|v\|^2+\|k\|^2)$ and the energy identity show
\[
 \min_{\substack{w\in\{v,k\}\\w\ne0}}
       \frac{q(w,w)}{\|w\|^2}
 \leq\frac{q(v,v)+q(k,k)}{\|v\|^2+\|k\|^2}
 \leq2\frac{q(u,u)}{\|u\|^2}.
\]
This construction changes no physical coefficient or state norm.
All assertions about sequences and uniform bounds follow directly
from the two-sided inequalities for each member of the sequence.
\end{proof}

This comparison calculates how a putative collective low-energy sector
must appear under the exact block map: it appears either among the
full zero-shift reconstructed states, with their eliminated norm
retained, or among physical states in the complementary sector.
The decomposition \eqref{eq:softdecomposition} exhibits the vectors
and controls their energy quotients by the numerical factor two.
It does not assert which alternative occurs in the interacting
large-volume limit, or that either one actually becomes soft.
The fixed-loop high-energy fraction in Theorem~\ref{thm:spectral}
continues to hold in its full scope. Equations
\eqref{eq:zeroloopactual}--\eqref{eq:zeroloopenergy} and
Theorem~\ref{thm:gaptransfer} now supply the proved zero-shift and
all-state relations without extending that fraction to an unsupported
family of vectors. For $m=1$, $J_1$ is the identity, $Q_1=B_1=0$,
and reconstruction is the identity; the physical gap relation is then
exactly $\lambda_1^{\rm R}=\delta_{\rm G}$ with no complementary sector.

\section{Physical skeleton morphism and uniform complementary inverse}
\label{sec:uniforminverse}

This section uses the complete fixed-point and physical-resolvent proof
retained in \nolinkurl{references/LANE_A_ELIMINATION_INPUTS.md},
Sections 1, 10, 11, 12 and 15. The input concerns the identical open
box, every original link and face, and the actual positive unit vacuum.
Here is its exact dictionary. Its $g_{\rm YM}$ is our $g$, its
$\mathcal E$ is our $E_0$, and its coordinate action
$\alpha_h(U)_e=h_{s(e)}U_eh_{t(e)}^{-1}$ obeys
\[
 (h\cdot U)_e=h_{s(e)}^{-1}U_eh_{t(e)}
                  =(\alpha_{h^{-1}}U)_e.
\]
Thus its Hilbert-space representation
$\mathsf T_hu=u\circ\alpha_{h^{-1}}$ is exactly
$\mathsf T_hu(U)=u(h\cdot U)$ here. The identity map on functions
intertwines these representations, on $L^2$, $H^1$ and $H^2$;
no assertion that parameter inversion is a group homomorphism is needed.
The two Hamiltonians and their domains coincide. Uniqueness of the
positive unit vacuum then proves equality of the two $\psi$, of
$\rho=\psi^2$, and of the ground energies. The source scalar in its
(12.7) will be written $c_{\rm vac}=\langle1,\psi\rangle$ here;
it is not the density constant $\alpha$ in \eqref{eq:finitegap}.

For precision, the proved input is
\begin{equation}\label{eq:stronggap}
 \delta_{\rm G}\geq\Delta_{\rm sc}
 :=3\kappa\left(1-\frac{512}{3}\xi\right)
 =\kappa(3-512\xi)\geq\frac{287\kappa}{96},
 \qquad 0<\xi\leq\frac1{49152}.
\end{equation}
Its entire proof is retained, not a proposed additional hypothesis.
The exact source construction maps a collection of nonconstant
Haar sectors $z_I$ to $\phi=e^{\mathfrak C}1$, where
$\mathfrak C=\sum_{I\ne\varnothing}|z_I\rangle\langle1_I|
\otimes I_{I^c}$ is bounded and nilpotent in each finite box.
The inverse is $e^{-\mathfrak C}$, with both exponentials finite
polynomials preserving $D(H_0)=H^2$. Its full fixed-point proof gives
\[
 \psi=c_{\rm vac}\phi,\quad
 c_{\rm vac}>0,\quad c_{\rm vac}^2\|\phi\|^2=1,\quad
 e^{-\mathfrak C}(H-E_0)e^{\mathfrak C}=H_0+\mathcal K,
 \qquad H_0=\kappa\sum_eE_e.
\]
These maps preserve the vacuum scalar and all of $2bM$ in the
eigenvalue equation. The source proves that every iterate, hence both
similarities, commutes with every vertex gauge transformation.
In the orthogonal sector decomposition a physical nonempty support
has no degree-one vertex: averaging at such a vertex would give both
the same invariant vector and its zero one-link Haar average.
Every nonempty surviving support therefore contains at least four
edges of this open bipartite cubic graph. With the original
$E_e$ Casimir, $H_0\geq3\kappa$ on the physical nonconstant sectors.
The complete interaction estimate proved there is
$\|\mathcal K f\|_{\oplus,1}\leq(512\xi/3)\|H_0f\|_{\oplus,1}$,
where $\|f\|_{\oplus,1}=\sum_I\|\Pi_If\|$ and
$\|f\|\leq\|f\|_{\oplus,1}\leq2^{N/2}\|f\|$.
The retained proof supplies every commutator and counting estimate
and the contraction on $\|z\|_*\leq1/16$, with constant at most
$11/36$ in the displayed range. The sector norm equivalence is
explicitly volume dependent. The physical Neumann resolvent excludes
every $0<t<3\kappa(1-512\xi/3)$, transfers under the displayed
domain-preserving similarity, and proves \eqref{eq:stronggap} for the
original self-adjoint operator. No form positivity of a nonunitary
similarity or deletion of its norm factors is used.

The retained proof derives numerical constants for the creation method
of Yarotsky~\cite{Yarotsky}, Section 2, source equations
\texttt{mmo}, \texttt{c1}, \texttt{vsum}, \texttt{thl},
\texttt{sumj}. Its physical-space reference is Baez~\cite{Baez},
the section ``Gauge Theory on a Graph,'' Lemmas 1, 2 and the source
lemma \texttt{lem2.5}. Baez's link convention
$A_e\mapsto h_{t(e)}A_eh_{s(e)}^{-1}$ is related to the input's
$U_e\mapsto h_{s(e)}U_eh_{t(e)}^{-1}$ by the involution
$\jmath(U)_e=U_e^{-1}$: direct multiplication gives
$\jmath(\alpha_hU)_e=h_{t(e)}\jmath(U)_eh_{s(e)}^{-1}$.
The pullback is a Haar-unitary map with itself as inverse,
preserves $H^1,H^2$ because inversion is a bi-invariant metric
isometry, and intertwines each $E_e$. The raw Wilson products in
our Hamiltonian are unchanged. Both original TeX sources are retained.

The source conditional complement retains individual fine links.
Our complement retains their ordered products. The following proves
the exact relation, including the Hilbert and operator domains.

\begin{proposition}\label{prop:physicalbridge}
Let $2\leq m\mid2L$, $s=(U_e)_{e\in S_m}$,
$r=(U_e)_{e\in R_m}$, and define
\[
 \mathcal A_m=\SU(2)^{R_m},\quad
 \bar\rho_R(r)=\int\rho(s,r)\dd\lambda_{S_m}(s),\quad
 \mathcal K_R=L^2(\mathcal A_m,\bar\rho_R\lambda_R).
\]
Let $J_S:\mathcal K_R\to\HH$ be $J_SF(s,r)=F(r)$, with
\[
 (J_S^*u)(r)=\frac{1}{\bar\rho_R(r)}
                \int\rho(s,r)u(s,r)\dd\lambda_{S_m}(s),
 \qquad P_S=J_SJ_S^*.
\]
Write $\mathcal K_R^{\rm G}$ for invariance under the endpoint action
of all original vertices. There is a unitary map of physical spaces
\begin{equation}\label{eq:skeletonunitary}
 \iota_m:\KK_m^{\rm G}\longrightarrow\mathcal K_R^{\rm G},
 \qquad (\iota_m f)(r)=f((u_{\gamma,1}\cdots u_{\gamma,m})_\gamma).
\end{equation}
Its inverse on arbitrary $L^2$ classes is
\begin{equation}\label{eq:skeletoninverse}
 (\iota_m^{-1}F)(z)=\int F(\Theta_m^{-1}(v,z))\dd\lambda_v(v),
 \quad
 \Theta_m:\mathcal A_m\longrightarrow
 \left(\prod_\gamma\SU(2)^{m-1}\right)\times\mathcal Q_m,
 \quad r\longmapsto(v,z).
\end{equation}
The inverse coordinates are $u_{\gamma,k}=v_{\gamma,k}$ for $k<m$
and $u_{\gamma,m}=(v_{\gamma,1}\cdots v_{\gamma,m-1})^{-1}z_\gamma$.
Both maps identify the invariant $H^1$ and $H^2$ spaces, and
\begin{equation}\label{eq:physicalprojections}
 J_S\iota_m=J_m,
 \qquad J_S^*|_{\HH^{\rm G}}=\iota_mJ_m^*|_{\HH^{\rm G}},
 \qquad P_S|_{\HH^{\rm G}}=P_m|_{\HH^{\rm G}}.
\end{equation}
Consequently the identity on functions is an isometry from the
source physical conditional complement for $S=S_m$ onto
$\HH_{Q_m}^{\rm G}$. It identifies their full form domains,
their complementary operator domains, their operators and their
resolvents. In the range \eqref{eq:stronggap},
\begin{equation}\label{eq:uniformC}
 C\geq\Delta_{\rm sc} I,\qquad
 \|C^{-1}\|\leq\Delta_{\rm sc}^{-1}
                  \leq\frac{96}{287\kappa}.
\end{equation}
\end{proposition}

\begin{proof}
The two inverse coordinate compositions and the Haar identity for
$\Theta_m$ follow by the exact prefix cancellations already proved
for $\Phi_m$, with all $S_m$ factors omitted from both sides of
this particular map. Integrating those factors in the density gives
\[
 \bar\rho_R(\Theta_m^{-1}(v,z))=r_m(z),\qquad
 \bar\rho_R\dd\lambda_R=r_m(z)\dd\lambda_v\dd\lambda_m.
\]
To prove the first equality, perform simultaneously the internal
vertex transformations with fixed coarse endpoints as in Section 3;
they send $(v,z)$ to $(1,z)$, preserve the marginal by substitution
on $S_m$, and leave $z$ unchanged. Integrating in $v$ then identifies
the remaining function with the definition of $r_m$.

There is also an $L^2$ proof of the descent, without evaluation on
a measure-zero section. For a chain with endpoints fixed, the
internal gauge action is
$v'_1=u_1h_1$ and $v'_k=h_{k-1}^{-1}u_kh_k$ for $1<k<m$.
For fixed input $r$, Haar integration successively over
$h_1,\ldots,h_{m-1}$ makes the variables $v'_1,\ldots,v'_{m-1}$
product Haar; the chain product remains $z$. Applying this to every
disjoint interior shows that the compact internal gauge average on
$L^2(\bar\rho_R\lambda_R)$ is precisely fibre integration in $v$
followed by constant lift. Hence every invariant $F$ equals that
lift almost everywhere. The remaining coarse endpoint action is
$z_\gamma\mapsto h_x^{-1}z_\gamma h_{x+m\mathbf e_i}$; this proves
the full coarse invariance of \eqref{eq:skeletoninverse} and the
fine invariance of \eqref{eq:skeletonunitary}. Substitution proves
both inverse identities. Moreover
\[
 \|\iota_m f\|_{\mathcal K_R}^2
 =\int |f(z)|^2r_m(z)\dd\lambda_v\dd\lambda_m
 =\|f\|_{\KK_m}^2.
\]
No vector is divided by its mass. Smooth coordinate pullback on
these compact products preserves each Sobolev order. Differentiating
under the compact fibre integral and applying Cauchy--Schwarz bounds
the inverse in $H^1$ and $H^2$. Smooth positive weights give the
same Sobolev sets in each finite box. Approximation, followed by
compact gauge averaging, proves the claimed domain identifications.

The first identity in \eqref{eq:physicalprojections} follows from
the coordinate product. For physical $u$, substitution in the
conditional integral shows that $J_S^*u$ is invariant under every
skeleton endpoint action. In $(v,z)$ coordinates it is therefore
independent of $v$ almost everywhere. Its $v$ integral equals
$r_m(z)^{-1}\int\rho(s,\Theta_m^{-1}(v,z))
u(s,\Theta_m^{-1}(v,z))\dd\lambda_{S_m}\dd\lambda_v$,
which is exactly $J_m^*u$. This proves the second identity and then
the third. It does not assert equality of the projections on all
noninvariant vectors.

Thus the two physical complements are the same closed subspace
of $\HH$, with the same form $q$, including derivatives in every
fine link. Both have form domain $H^1$ intersected with that space.
The weak operator-domain rule is, in both presentations,
\[
\begin{aligned}
 D(C)=\bigl\{u\in H^1\cap\HH_{Q_m}^{\rm G}:{}&
 \text{ there exists }w\in\HH_{Q_m}^{\rm G}\text{ such that}\\
 &q(v,u)=\langle v,w\rangle
 \text{ for every }v\in H^1\cap\HH_{Q_m}^{\rm G}\bigr\},\\
 Cu={}&w.
\end{aligned}
\]
Equality of the tests, inner products and forms proves equality
of the operators and their resolvents, not merely of their formal
differential expressions. Finally every vector in the complement
has $\mu u=0$, and $\mathcal Uu=\psi u$ is physical and perpendicular
to $\psi$. The already-proved \eqref{eq:stronggap} applied under
this unchanged unitary gives \eqref{eq:uniformC}.
\end{proof}

\begin{theorem}\label{thm:uniformzero}
In the range \eqref{eq:stronggap}, for every mean-zero
$f\in\VV_m\cap\KK_m^{\rm G}$, put $b_f=B_mf$,
$d_f=\|f\|^2$, $a_f=a_m(f,f)$ and retain the raw measure
$\sigma_f(I)=\|1_I(C)b_f\|^2$.
Then
\begin{equation}\label{eq:uniformzeroenergy}
 s_{m,0}(f,f)\geq\Delta_{\rm sc}
              (d_f+\|C^{-1}b_f\|^2),\qquad
 \lambda_m^{\rm R}\geq\Delta_{\rm sc},\quad c\geq\Delta_{\rm sc}.
\end{equation}
For the unchanged mean-centered simple loop $f_C$, let
$\ell$ be its fine perimeter and $d_C,r_C$ retain their earlier
definitions. The physical zero-shift state $u_0=\psi\mathcal R_m f_C$
has the exact mass and energy
\begin{align}
 \|u_0\|^2&=d_C+\int t^{-2}\dd\sigma_{f_C}(t),\label{eq:uniformloopmass}\\
 \mathfrak q_{H-E_0}[u_0]
 &=\kappa\ell(1-r_C/4)-\int t^{-1}\dd\sigma_{f_C}(t)
 \geq\Delta_{\rm sc}\|u_0\|^2.\label{eq:uniformloopenergy}
\end{align}
Its quantitative upper mass bound is
\begin{equation}\label{eq:uniformloopbound}
 \|u_0\|^2\leq d_C+
 \min\left\{\frac{\kappa\ell(1-r_C/4)}{\Delta_{\rm sc}},
             \frac{16G^2\kappa^2\ell^2}{\Delta_{\rm sc}^2}\right\}
 \leq4+\frac{96\ell}{287}.
\end{equation}
For every $z\geq0$, define $w_z=(C+z)^{-1}b_f$,
$N_f(z)=d_f+\|w_z\|^2$ and
$e_f(z)=a_f-\int(t+z)^{-1}\dd\sigma_f
                    -z\int(t+z)^{-2}\dd\sigma_f$.
These are the full physical mass and excitation energy of
$\psi(J_mf-w_z)$. The exact differences are
\begin{align}
 w_0-w_z&=zC^{-1}(C+z)^{-1}b_f,\label{eq:uniformwdiff}\\
 e_f(z)-e_f(0)&=\int\frac{z^2}{t(t+z)^2}\dd\sigma_f(t),
                                                        \label{eq:uniformediff}\\
 N_f(0)-N_f(z)&=\int\left(\frac1{t^2}-\frac1{(t+z)^2}\right)
                                      \dd\sigma_f(t),    \label{eq:uniformndiff}\\
 s_{m,z}(f,f)-s_{m,0}(f,f)-zd_f
 &=\int\frac{z}{t(t+z)}\dd\sigma_f(t).    \label{eq:uniformmemorydiff}
\end{align}
Writing $\theta=z/(\Delta_{\rm sc}+z)$, their estimates are
\begin{align}
 \|w_0-w_z\|&\leq\theta\sqrt{a_f/\Delta_{\rm sc}},\nonumber\\
 0\leq e_f(z)-e_f(0)&\leq\theta^2a_f,\nonumber\\
 0\leq N_f(0)-N_f(z)&\leq
 \left[1-\left(\frac{\Delta_{\rm sc}}{\Delta_{\rm sc}+z}\right)^2\right]
 \frac{a_f}{\Delta_{\rm sc}},\nonumber\\
 0\leq s_{m,z}(f,f)-s_{m,0}(f,f)-zd_f&\leq\theta a_f.
                                                        \label{eq:uniformcontinuity}
\end{align}
The constants controlling the inverse do not depend on box or block
size. The displayed $a_f$, $\ell$, $G$, $\kappa$, raw norms and
spectral measures keep their full dependence.
\end{theorem}

\begin{proof}
The map $\mathcal R_m$ and its exact domain and energy identities
were proved in Proposition~\ref{prop:finitezero}. Its image under
$\mathcal U$ is physical and perpendicular to $\psi$. Applying
\eqref{eq:stronggap} proves \eqref{eq:uniformzeroenergy}, including
the infimum inequality. The previous exact inverse-moment formulas
give \eqref{eq:uniformloopmass}--\eqref{eq:uniformloopenergy}.
Positivity of $s_{m,0}$ gives
$\int t^{-1}\dd\sigma_f\leq a_f$. Since
$\operatorname{supp}\sigma_f\subset[\Delta_{\rm sc},\infty)$,
\[
 \int t^{-2}\dd\sigma_f\leq
 \min\{a_f/\Delta_{\rm sc},\|b_f\|^2/\Delta_{\rm sc}^2\}.
\]
Insert the original loop identities $a_f=\kappa\ell(1-r_C/4)$,
$\|b_f\|\leq4G\kappa\ell$, $d_C\leq4$, $a_f\leq\kappa\ell$
and \eqref{eq:stronggap} to prove \eqref{eq:uniformloopbound}.

The resolvent identity proves \eqref{eq:uniformwdiff}. Subtraction
of the complete energy integrands gives
$t^{-1}-(t+z)^{-1}-z(t+z)^{-2}=z^2/[t(t+z)^2]$.
The norm subtraction and the definition of $s_{m,z}$ give the
remaining exact differences. On the support of the raw measure,
$z/(t+z)\leq\theta$ and
$1-[t/(t+z)]^2\leq1-[\Delta_{\rm sc}/(\Delta_{\rm sc}+z)]^2$.
Apply these inequalities to the corresponding nonnegative
integrands and use the preceding inverse-moment bounds. Squaring
\eqref{eq:uniformwdiff} first proves its norm estimate. This proves
every inequality in \eqref{eq:uniformcontinuity}, with no removal
of the measure or its total mass $\|B_mf\|^2$.
\end{proof}

For an elementary coarse face $\ell=4m$, the last mass bound is
$4+384m/287$, and its physical side length remains $ma$.
In original parameters the new lower energy scale and range are
\[
 \Delta_{\rm sc}=\frac{2g^2}{a}\left(3-\frac{128}{g^4}\right)
 \geq\frac{287g^2}{48a},\qquad g^4\geq12288.
\]
The expectation of $H$ is still $E_0N_f(z)+e_f(z)$.
This is a statement for the interacting finite boxes, uniform as
their sizes vary within this explicit coupling range. It takes no
spatial continuum limit and provides no small-$g$ continuation.
Every induced score covariance and the full resolvent and time memory
remain those of Sections 4--6.


\section{Original scales and the exact scope of the spectral conclusion}

For an elementary coarse face $\ell=4m$. The actual state has
\begin{align}
 \|u_C\|^2&=r_C-(\overline W_C)^2\in[\alpha,4],\label{eq:physicalnorm}\\
 \mathfrak q_{H-E_0}[u_C]
 &=\frac{8g^2m}{a}\left(1-\frac{r_C}{4}\right)
       \in\left[\frac{6\alpha g^2m}{a},\frac{8g^2m}{a}\right],\label{eq:physicalenergy}\\
 \frac{\mathfrak q_{H-E_0}[u_C]}{\|u_C\|^2}
 &\in\left[\frac{3\alpha g^2m}{2a},\frac{8g^2m}{a\alpha}\right],
 \qquad \Omega_C=\frac{3\alpha g^2m}{4a}.\label{eq:physicalquotient}
\end{align}
Here
$\alpha=1/(9\,12^{16/g^4})$, and $G,C_*,c_H$ are the explicit functions
\eqref{eq:G}, \eqref{eq:highweight} of $d=4/g^4$.
The physical length is $ma$, so $g^2m/a=g^2(ma)/a^2$ as an exact
rewriting with no coordinate change. The expectation of the unsubtracted
Hamiltonian is
\[
 \langle u_C,Hu_C\rangle
 =E_0(r_C-(\overline W_C)^2)
       +\frac{8g^2m}{a}(1-r_C/4).
\]
The same scalar $E_0\|u_{z_C}\|^2$ is added to the dressed excitation
energy. Also $0\leq E_0\leq2bM$ follows from $H\geq0$ and the constant
unit Haar trial vector, whose plaquette means vanish. The original
$2bM$ term has not been omitted from $H$.

The exact instantaneous kinetic term has the two metric expressions
\[
 -\kappa m\Delta^T=-(\kappa m/4)\Delta^c,\qquad \kappa m=2g^2m/a.
\]
If its coefficient is written solely as
$2g_m^2/(ma)$, algebra forces $g_m^2=m^2g^2$. This is only the
coefficient dictionary for this marginal kinetic form. It does not make
\eqref{eq:marginalpotential} and the energy-dependent memory into the
original one-plaquette Hamiltonian at a claimed running coupling.

The map from the spatial coarse observable to the quantum theory is
exactly $f\mapsto\psi J_mf$, followed by \eqref{eq:nu}; the retained
resolvent is \eqref{eq:resolvent}, and sequential eliminations obey
\eqref{eq:recursive}. Theorems~\ref{thm:spectral} and
Corollary~\ref{cor:dressedweight} are statements about that full quantum
spectral measure, with no omission of generated interactions. For fixed
$g>0$, these particular large-loop channels retain a positive fraction
of weight at energies proportional to $g^2m/a$, including the specified
fully dressed channels. Their unnormalized spectral mass above the
displayed threshold is bounded below by the stated positive multiple
of their unchanged squared norm, at every increasing $m$ and fixed $a,g$.

This assertion does not give a lower bound on the support of the whole
measure: some weight can lie below $\Omega_C$. Nor does a lower bound
on the quotient of this family bound the infimum over all physical
states. At variable $g=g(a)$ the constants $\alpha$ and $c_H$ have the
explicit coupling dependence above and can tend to zero. The calculation
therefore proves neither a continuum mass gap nor a four-dimensional
mass-gap counterexample. It supplies an exact extensive block map and
an all-positive-coupling, volume-independent quantitative channel result;
Theorem~\ref{thm:gaptransfer} additionally relates every physical
low-energy state to the zero-shift reconstructed and complementary
sectors, with parameter-independent comparison constants and explicit
state maps. Section~\ref{sec:uniforminverse} now supplies the proved
physical lower bound $\Delta_{\rm sc}=\kappa(3-512\xi)$, uniform
in box and block size for $0<\xi\leq1/49152$, together with the
exact skeleton/complement morphism and full-memory zero-shift bounds.
The all-positive-coupling finite bound \eqref{eq:finitegap} remains
valid with its displayed volume dependence. Neither the comparison
nor the exhaustive affine calculation
in Theorem~\ref{thm:affine} presumes which sector becomes soft or an
answer for the joint spatial continuum limit.

\section{Primary-source map and reproducibility}

The general Feshbach--Schur construction is attributed to its established
literature. The locally retained source of Dusson, Sigal and Stamm was
read at Theorem~1.2, its reconstruction map, and Eqs.~(1.10)--(1.16).
Their Corollary~1.3 is the energy-dependent fixed-point formulation.
Its full eigenvalue condition tests all retained directions; the
scalar affine root \eqref{eq:affineroot} tests the fixed retained
direction only, as proved after Theorem~\ref{thm:affine}.
Their operator $H-\lambda$ corresponds here to $\LL+z$, with
$\lambda=-z$, their complementary inverse to $(C_m+z)^{-1}$, and
their reconstruction to \eqref{eq:reconstruction}. Their effective
interaction has the negative sign in \eqref{eq:schur}. The infinite-rank
form argument and the conditional coefficients are proved in this paper.

Bakry and Ledoux's reverse Poincare estimate is their Eq.~(1.7), with
$d(t)=2t$ at curvature bound zero, on pp.~686--687. Our single-edge
Markov generator is $\kappa\Delta^T$ and its carré du champ is
$\kappa\sum_a(X_af)^2$, giving precisely \eqref{eq:reverse}. Its full
compact-group proof was included. Simon's locally retained Theorem~1.1
uses the Euclidean generator $-\Delta/2$ and its Brownian-bridge formula.
It is a primary reference for positive potential weights, not a theorem
identifying Euclidean space with this compact configuration product;
the compact parabolic comparison needed here was proved in
\eqref{eq:comparison}. Mori's projection-memory method provides the
historical comparison: its projection is Eq.~(2.13) and its exact
memory equations are Eqs.~(3.1)--(3.10), pp.~430--431.
Our evolution uses $-\LL$ on functions with the vacuum inner product;
Mori's evolution uses a Liouville operator on dynamical variables.
Equation~\eqref{eq:memory} derives the corresponding Euclidean block
identity here, including its signs and domains.

The companion exact checker verifies the entire open-box skeleton and
path counts on several finite boxes, chain disjointness, nested product
paths, every retained loop orientation, exact Pauli rotations and the
mixed derivative tensor in explicit noncommuting samples, and exact
rational Schur associativity and spectral constants.
The zero-shift checker separately verifies physical norm and quotient
derivatives, all six interior/boundary and coupled/uncoupled affine
examples, the complete energy and norm decomposition, and the harmonic
gap comparison in 81 rational interacting block examples.
The physical-complement checker additionally verifies noncommuting
chain gauge actions, their inverse compositions, the exact original
strong-coupling constants, and the full inverse-moment continuity
identities and inequalities. These finite checks
diagnose algebra; they do not stand in for the all-box proofs above or
compute a substitute vacuum. No Lean run or large parallel job is used.

\section{Material-tensor transfer into interacting local-energy states}
\label{sec:fabel-tensor}
This section retains the original Wilson Hamiltonian, vacuum, gauge action,
spacing and coupling from Section~1. It constructs a particular observable
map from the full Fabel material tensor. No equality between fluid evolution
and quantum Hamiltonian evolution is asserted. The map and its exact
spectral measures are calculated below; the remaining infinite-volume
spectral question is not replaced by the value of a coordinate.

\subsection{Transport conventions and the vacuum-moment correction}
For the column-section connection $d+\mathcal A$, ordinary parallel transport
along an oriented edge $e:[0,a]\to\mathbb R^3$ solves
\[
 P_e'(s)=-\mathcal A(e(s))(\dot e(s))P_e(s),\qquad P_e(0)=I.
\]
Its gauge transform is $P_e^h(s)=h(e(s))^{-1}P_e(s)h(e(0))$.
Differentiating this formula verifies the transformed ODE and its initial
value; uniqueness gives the assertion. The links in Section~1 have the
opposite endpoint convention and are therefore \emph{inverse} transports:
\begin{equation}\label{eq:fabel-inverse-transport}
 U_e=P_e(a)^{-1},\quad
 U_e^h=h(e(0))^{-1}U_eh(e(a)),\quad U'(s)=U(s)\mathcal A(e(s))(\dot e(s)).
\end{equation}
Inverses and concatenations follow by multiplication of the corresponding
transport maps. Thus this is an exact gauge-equivariant map for arbitrary
noncommuting connection components. It is not a global inverse from finite
links to connections.

The Cartan specialization $\mathcal A_i=\lambda u_iT$, $T=-i\sigma_3/2$,
has
\[
 U_e=\exp\left(\lambda T\int_e u\cdot dx\right),\qquad
 U_p=\exp(\vartheta_pT),\quad
 \vartheta_p=\lambda\oint_{\partial p}u\cdot dx .
\]
With positively oriented $(i,j)$ faces, Stokes' theorem gives
$\vartheta_p=\lambda\int_p(\partial_i u_j-\partial_j u_i)\,dx^i dx^j$.
For a fixed smooth field on a compact region it equals
$a^2\lambda(\partial_i u_j-\partial_j u_i)+O(a^3)$.
In the general non-Abelian case the four edge expansions to second order
give $U_p=I+a^2F_{ij}+O(a^3)$: the differentiated terms are
$\partial_i\mathcal A_j-\partial_j\mathcal A_i$ and the four products
give $[\mathcal A_i,\mathcal A_j]$, with this positive sign.
For uniform $C^2$ bounds the remainder is uniform on that compact region.
As $\tr e^{\vartheta T}=2\cos(\vartheta/2)$, the exact magnetic value is
\begin{equation}\label{eq:fabel-magnetic-value}
 b(2-W_p)=\frac{2-2\cos(\vartheta_p/2)}{2g^2a}
       =\frac{\vartheta_p^2}{8g^2a}
          +O_g(\vartheta_p^4/a).
\end{equation}
For fixed $g,\lambda$ and a fixed smooth field, summation on a regular
mesh of a bounded region yields
$\lambda^2(8g^2)^{-1}\int|\operatorname{curl}u|^2dx$.
This follows by a Riemann sum; each leading cell term is $a^3$ times the
displayed density and each uniform cell error is $O(a^4)$.
Pointwise curvature growth alone does not imply growth of this integral:
the volume of the concentrating region must also be accounted for.

In Lorentzian coordinates $X^0=ct,X^i=x^i$, the same Cartan field with
$\mathcal A_0=0$ has the complete conserved current
\begin{align*}
 j_0&=0,\\
 j_i&=\frac{\lambda}{g^2}
       (\Delta u_i-c^{-2}\partial_t^2u_i)T\\
 &=\frac{\lambda}{g^2}\left[
 \Delta u_i-c^{-2}\left(\nu\Delta w_i
 -\sum_jw_j\partial_ju_i-\sum_ju_j\partial_jw_i
 -\partial_i\partial_tp+\partial_tf_i\right)\right]T,\qquad w_i=\partial_tu_i.
\end{align*}
This is obtained by differentiating the original Navier--Stokes equation,
not by changing its viscosity or force. $D^\nu j_\nu=0$ follows from
incompressibility and the common Cartan factor. The field solves that
sourced classical equation. It can be used as data for a quantum trial
state without claiming it solves source-free Yang--Mills. Proving
source-free classical dynamics is not a prerequisite for defining a trial
state; its physical Hilbert and spectral properties must instead be proved.

There is a separate important correction concerning Wilson observables.
Evaluating $W_C$ at the links of a background gives a number $W_C(U[\mathcal A])$.
The functional $W_C(U)$ on the full configuration space has no dependence
on the chosen background until an additional prescription supplies it.
The generic centered state $\psi(W_C-\mu W_C)$ is physical and nonzero,
but naming a background does not make it a background-dependent state map.
Moreover $r_C=\mu(W_C^2)$ is the actual quantum second moment. It is not
$W_C(U[\mathcal A])^2$.

Keep the exact bound from Section~7,
$\alpha\le r_C\le4-3\alpha$, with
$\alpha=(9\,12^{16/g^4})^{-1}>0$. Then
\begin{equation}\label{eq:fabel-loop-scale-corrected}
 \frac{3\alpha}{4}\kappa\ell\le
 a_C=\kappa\ell(1-r_C/4)
 \le(1-\alpha/4)\kappa\ell .
\end{equation}
For physical perimeter $P_C=a\ell$ this bounds the \emph{bare} loop
numerator between the displayed constants times $2g^2P_C/a^2$ at fixed $g$;
it does not give asymptotic ratio one. The reconstructed numerator and norm
remain $a_C-M_1(0)$ and $d_C+M_2(0)$.
The lower bound $\Delta_{\rm sc}$ for these reconstructed states is used
only for $\xi\le1/49152$. In a spatial continuum path $a\downarrow0$,
$\kappa\to0$ would imply $g^2=(a/2)\kappa\to0$ and thus $\xi\to\infty$.
Without $a\downarrow0$ the latter implication is not valid.

\subsection{The full weighted local-energy operator}
For arbitrary real weights $f=(f_e)_{e\in\mathsf E_L}$ define
\begin{equation}\label{eq:fabel-local-D}
 f_p=\frac14\sum_{e\in\partial p}f_e,\quad
 V_p=b(2-W_p),\quad
 D_f=\kappa\sum_ef_eE_e+\sum_pf_pV_p,\quad
 \Xi_f=(D_f-\langle\psi,D_f\psi\rangle)\psi .
\end{equation}
All sums include the boundary edges and faces.
The electric-only operator $\sum f_eE_e$ is a different operator.
On a Peter--Weyl product sector with edge spins $j_e$, the signed electric
sum in $D_f$ has eigenvalue $\kappa\sum_ef_ej_e(j_e+1)$.
Its domain consists of square-summable coefficients with the squares of
these eigenvalues as weights. Addition of the bounded real magnetic
multiplication preserves that self-adjoint domain. Smooth functions lie in
the domain, and $D_f$ maps smooth functions to smooth functions.
In particular $\Xi_f$ is smooth and belongs to every power domain of $H$.
Gauge invariance of each link Casimir and each face trace gives
gauge invariance of $\Xi_f$. Its definition proves
$\langle\psi,\Xi_f\rangle=0$ with no rescaling.

Put $\mathcal E=E_0$ and $\mathcal A=H-\mathcal E$.
The product rule, including both derivative terms, gives
\begin{align}
 [E_e,V_p]F&=-\sum_a\bigl[(X_{e,a}^2V_p)F+
                           2(X_{e,a}V_p)X_{e,a}F\bigr],\label{eq:fabel-local-comm}\\
 [H,D_f]&=\kappa\sum_p\sum_{e\in\partial p}(f_p-f_e)[E_e,V_p].
 \label{eq:fabel-current-full}
\end{align}
The first expression vanishes when $e\notin\partial p$.
The second identity follows by expanding the commutator: electric terms
commute with electric terms, magnetic multiplications commute, and the
two remaining electric--magnetic sums have the stated coefficient.
It proves
\[
 \mathcal A\Xi_f=[H,D_f]\psi,\qquad
 q_{\mathcal A}[\Xi_f]=\tfrac12
 \langle\psi,[D_f,[H,D_f]]\psi\rangle .
\]
For the raw positive measure
$\nu_f(B)=\|1_B(\mathcal A)\Xi_f\|^2$, $B\subset(0,\infty)$, spectral
integration yields
\begin{equation}\label{eq:fabel-three-moments}
 \|\Xi_f\|^2=\int d\nu_f,\quad
 q_{\mathcal A}[\Xi_f]=\int E\,d\nu_f,\quad
 \|[H,D_f]\psi\|^2=\int E^2d\nu_f .
\end{equation}
All moments are finite by smoothness. For a constant $f_e=c_0$,
$D_f=c_0H$, so its expectation is $c_0\mathcal E$ and $\Xi_f=0$.
This is an exact cancellation, not a zero-energy nonzero state.

\subsection{Coordinate-level cubic symmetries and the full covector map}
Write $m(e)=n+\mathbf e_i/2$ for edge midpoints. A signed permutation
matrix $R$ maps the vertex box to itself. If $R\mathbf e_i=\epsilon\mathbf e_j$,
then the positive image edge of $e=(n,i)$ is $(Rn,j)$ for $\epsilon=1$,
and $(Rn-\mathbf e_j,j)$ for $\epsilon=-1$.
Define $\Gamma_R(U)$ at this image edge to be $U_e^\epsilon$.
The inverse is $\Gamma_{R^{-1}}$; composing the signs and endpoints
verifies both inverse compositions. Every midpoint is mapped to $Rm(e)$.
Each face becomes a face, possibly with its orientation reversed.
Its SU(2) trace is unchanged under reversal.
Inversion of a group variable is Haar preserving and takes the
left Casimir to the equal right Casimir, since the metric is bi-invariant.
Thus the pullback unitary
$\mathsf T_RF=F\circ\Gamma_R^{-1}$ preserves $H$, $H^1$ and $H^2$ and
intertwines the full gauge actions by $h'(Rn)=h(n)$.
Uniqueness of the positive unit vacuum gives $\mathsf T_R\psi=\psi$.
For weights transported by $f_R(e')=f(e)$ it follows that
$\mathsf T_RD_f\mathsf T_R^{-1}=D_{f_R}$ and
$\mathsf T_R\Xi_f=\Xi_{f_R}$. Spectral projections of $\mathcal A$ commute
with $\mathsf T_R$.

Let $\Xi_r=\Xi_{m_r}$. Reflection in coordinate $r$ changes $\Xi_r$ to
$-\Xi_r$ and fixes $\Xi_s$ for $s\ne r$; permutations exchange the three.
Consequently, for every real $t\in\mathbb R^3$ and $c_0\in\mathbb R$,
\begin{equation}\label{eq:fabel-affine-raw}
 \Xi_{c_0+t\cdot m}=\sum_rt_r\Xi_r,\quad
 \langle D_{c_0+t\cdot m}\rangle=c_0\mathcal E,\quad
 \nu_{c_0+t\cdot m}(B)=|t|^2\nu_{m_1}(B).
\end{equation}
Indeed a reflection negates each off-diagonal matrix element
$\langle\Xi_r,1_B(\mathcal A)\Xi_s\rangle$ while leaving it invariant;
it is therefore zero. Permutations equate the diagonal elements.
This proof applies at every positive coupling, including when a state
or a measure vanishes.

For the retained Fabel time parameter $z=\sqrt{1-8\tau}>0$,
$\tau=(1-z^2)/8$, the full inverse material derivative is
computed from the unchanged real polynomial map
\begin{align*}
 F_1(x,y,w)&=(1+xy)^3w+y^2(1+xy)(4+3xy),\\
 F_2(x,y,w)&=y+3x(1+xy)^2w+3xy^2(4+3xy),\\
 F_3(x,y,w)&=2x-3x^2y-x^3w .
\end{align*}
Here $w$ denotes the original third spatial coordinate, not time.
Expanding the three-by-three determinant of its differentiated matrix
gives $\det DF=-2$. At the retained points the full matrices are
\[
 J_0=\begin{pmatrix}-9/16&-3/8&-1/8\\3/8&25/4&3/4\\-17/2&-3&-1\end{pmatrix},
 \quad
 J_\gamma=\begin{pmatrix}
 -9z^3/16&-3z/8&-1/8\\3z^2/8&25/4&3/(4z)\\-17/2&-3/z^2&-1/z^3
 \end{pmatrix}.
\]
Direct matrix multiplication gives
\begin{equation}\label{eq:fabel-C-full}
 C_\tau=
 \begin{pmatrix}
 (17-9z^3)/8&(3-3z^3)/(4z^2)&(1-z^3)/(4z^3)\\
 (51-24z^2-27z^3)/16&(9+8z^2-9z^3)/(8z^2)&(3-3z^3)/(8z^3)\\
 (-153+36z^2+117z^3)/8&(-27-12z^2+39z^3)/(4z^2)&(-9+13z^3)/(4z^3)
 \end{pmatrix}.
\end{equation}
It is $J_0^{-1}B_\tau^{-1}J_\gamma$, where $J=DF$ for the original
polynomials and
\[
q_0=(1,-3/2,13/2),\quad
\gamma=(z^{-1},-3z/2,13z^2/2),\quad B_\tau=I+6\tau E_{23}.
\]
The factorization gives the inverse $A_\tau=J_\gamma^{-1}B_\tau J_0$
and determinant one.
The source matrix entries are reproduced and independently rederived
by the companion checker from $F$.

The covector $\zeta_\tau=C_\tau^{\mathsf T}e_3$ has the three entries
in the last row of \eqref{eq:fabel-C-full}. Choose the explicit affine
weight $f_\tau(e)=\zeta_\tau\cdot m(e)$. Then
\[
 \nu_{f_\tau}(B)=Q(z)\nu_{m_1}(B),
\]
where the entire, unmodified scalar is
\begin{align*}
 Q(z)={}&\frac{(-153+36z^2+117z^3)^2}{64}
 +\frac{(-27-12z^2+39z^3)^2}{16z^4}
 +\frac{(-9+13z^3)^2}{16z^6},\\
 &\lim_{z\downarrow0}z^6Q(z)=81/16 .
\end{align*}
Thus the full covector, with its changing direction, is realized by this
specified parameter-to-state map. Its raw norm, energy, and second spectral
moment all have the same factor $Q(z)$; their ratio does not depend on
$\tau$ whenever the state is nonzero. This is not an identification of
the Fabel cotangent space with the whole quantum state space.

\subsection{A non-affine map using every entry of the material tensor}
Set $K_\tau=C_\tau C_\tau^{\mathsf T}$ and keep all six independent entries.
For any real symmetric matrix $S$ define
\begin{equation}\label{eq:fabel-quadratic-map}
 f_S(e)=m(e)^{\mathsf T}S\,m(e),\qquad
 \mathfrak M_{L,a,g}(S)=\Xi_{f_S}\in
 C^\infty(\mathcal Q)^{\mathcal G}\cap\{\psi\}^{\perp}.
\end{equation}
This is a linear, typed map on the six-dimensional real matrix space.
It does not change the Hamiltonian metric: $S$ is the spatial weight of
the observable $D_f$, not a replacement kinetic tensor in $H$.
Using physical midpoints $a\,m(e)$ multiplies $f_S$ and $\Xi_{f_S}$ by
$a^2$, and every raw spectral measure by $a^4$.
Those factors remain present if physical midpoint units are chosen.

Use the three unchanged seed functions
\[
 f_0=m_1^2+m_2^2+m_3^2,\quad
 f_{\rm d}=m_1^2-m_2^2,\quad f_{\rm o}=2m_1m_2,
\]
and write $\nu_0,\nu_{\rm d},\nu_{\rm o}$ for their actual centered-state
spectral measures. These are positive measures of the full interacting
Hamiltonian, not free-channel approximations.

\begin{theorem}\label{thm:fabel-quadratic-measure}
For every fixed $L\ge2$, $a,g>0$, every real symmetric $S$ and every Borel
$B\subset(0,\infty)$,
\begin{align}
 \nu_{f_S}(B)&=w_0(S)\nu_0(B)+w_{\rm d}(S)\nu_{\rm d}(B)
                         +w_{\rm o}(S)\nu_{\rm o}(B),\label{eq:fabel-sector-measure}\\
 w_0(S)&=(\tr S)^2/9,\nonumber\\
 w_{\rm d}(S)&=\tfrac12\sum_{i=1}^3(S_{ii}-\tr S/3)^2,\quad
 w_{\rm o}(S)=\sum_{i<j}S_{ij}^2.\nonumber
\end{align}
The same identity holds after integration against $1,E,E^2$ and
$e^{-sE}$ for $s\ge0$, retaining the corresponding raw norms,
excitation energies, current norms, and time correlations.
\end{theorem}
\begin{proof}
Put $u_i=\Xi_{m_i^2}$ and $v_{ij}=\Xi_{2m_im_j}$.
A coordinate reflection fixes each $u_i$ and negates $v_{ij}$ if it
reflects just one of $i,j$. Thus every spectral inner product between
a $u$ and a $v$ vanishes. For two different off-diagonal pairs there
is a reflection negating exactly one, so their cross term also vanishes.
Permutations give the same diagonal spectral norm, equal to $\nu_{\rm o}(B)$,
for the three $v_{ij}$.

For the three $u_i$, permutations equate diagonal elements, denoted
$d(B)$, and equate off-diagonal elements, denoted $h(B)$.
A permutation exchanging two indices shows that $h(B)$ equals its
complex conjugate; it is real. The diagonal contribution for $s_i=S_{ii}$
is
\[
 d\sum_i s_i^2+2h\sum_{i<j}s_is_j
 =(d-h)\sum_i(s_i-\bar s)^2+(d+2h)(\tr S)^2/3,
 \qquad\bar s=\tr S/3.
\]
But $\nu_0=3d+6h$ and $\nu_{\rm d}=2(d-h)$ by the definitions of the
unchanged seed states. Substitution gives \eqref{eq:fabel-sector-measure}.
All identities are identities of finite positive measures. Integration
against the listed nonnegative functions is justified by monotone
convergence; smoothness makes the moments finite.
\end{proof}

Let $\eta=(1/4,3/8,-9/4)^{\mathsf T}$.
Every entry of \eqref{eq:fabel-C-full} yields
$z^3C_\tau\to\eta e_3^{\mathsf T}$ and hence
\[
 z^6K_\tau\longrightarrow S_*=\eta\eta^{\mathsf T}
 =\frac1{64}\begin{pmatrix}4&6&-36\\6&9&-54\\-36&-54&324\end{pmatrix}.
\]
Consequently the full nonlinear spatial profile
$m^{\mathsf T}K_\tau m$ has the following exact endpoint:
\begin{equation}\label{eq:fabel-tensor-limit-measure}
 z^{12}\nu_{f_{K_\tau}}\longrightarrow
 \frac{113569}{36864}\nu_0+
 \frac{100825}{12288}\nu_{\rm d}+
 \frac{531}{512}\nu_{\rm o}.
\end{equation}
This is convergence in total variation at each fixed regulator and
convergence of the first two moments. To prove it, note that each
$w_j$ is homogeneous of degree two and continuous in $S$, and apply
\eqref{eq:fabel-sector-measure}. The total variation of the difference
is bounded by the sum of the three coefficient errors times the finite
masses of the seed measures; the moment versions use their finite moments.
The constants follow by substituting $S_{*,ii}=4/64,9/64,324/64$ and
$S_{*,12}=6/64,S_{*,13}=-36/64,S_{*,23}=-54/64$.
Equivalently $z^6\Xi_{f_{K_\tau}}\to\Xi_{f_{S_*}}$ in every Sobolev norm
at a fixed regulator, by the linear six-state expression.
These displayed rescaled limits record the growth of the raw state;
the state itself has not been divided by its norm.

All three coefficients in \eqref{eq:fabel-tensor-limit-measure} are positive.
For sufficiently small $z>0$, each $w_j(K_\tau)$ is positive as well.
Then the positive-energy support of $\nu_{f_{K_\tau}}$ equals the union
of the supports of the nonzero seed measures: a Borel neighborhood
has zero mass in the sum precisely when it has zero mass in each seed.
Write $d_j=\int d\nu_j$, $n_j=\int E\,d\nu_j$.
If at least one $d_j$ is nonzero, the limiting quotient is exactly
\begin{equation}\label{eq:fabel-tensor-limit-quotient}
 \lim_{\tau\uparrow1/8}
 \frac{q_{\mathcal A}[\Xi_{f_{K_\tau}}]}{\|\Xi_{f_{K_\tau}}\|^2}
 =
 \frac{\frac{113569}{36864}n_0+
       \frac{100825}{12288}n_{\rm d}+\frac{531}{512}n_{\rm o}}
      {\frac{113569}{36864}d_0+
       \frac{100825}{12288}d_{\rm d}+\frac{531}{512}d_{\rm o}} .
\end{equation}
This is an evaluated relation between actual moments, not an estimate
that any unknown moment is small. If all $d_j=0$, every state in this
quadratic family is zero. The next calculation proves that this latter
case does not occur for sufficiently small positive $\xi$ on any fixed box.

\subsection{The first surviving state and exact open-boundary coefficients}
The full local-energy source proof is retained unchanged as
\path{references/LOCAL_ENERGY_CURRENT_INPUT.md}. We reproduce the
coefficient argument needed here. Set $H_0=\sum_eE_e$ and
$\mathcal W=\sum_pW_p$. On a fixed box, $\|\mathcal W\|_\infty=2M$ and
the nonconstant free gap is $3/4$.
The circle $|\zeta|=3/8$ has free resolvent norm at most $8/3$, so the
Neumann series for $H_0-\xi\mathcal W$ converges for $|\xi|<3/(16M)$.
The simple ground projection persists there. For real $\xi$ near zero
its positive unit vacuum is real analytic in each Sobolev space, by the
eigen-equation and elliptic regularity. Keep its unit-vacuum scalar:
\begin{align*}
 \psi_\xi&=1+\xi\mathcal W/3+\xi^2 B+O_{H^k,L}(\xi^3),\\
 H_0B&=(\mathcal W^2-M)/3,\qquad \int B=-M/18.
\end{align*}
For two faces sharing one edge, decompose
$W_pW_q=P_{pq,0}+P_{pq,1}$ by Haar averaging that edge.
The first term is one half the six-edge boundary trace, by the
fundamental matrix-element identity $\int U_{ij}\overline U_{kl}
=\delta_{ik}\delta_{jl}/2$.
The shared-edge spins are zero and one, and all six outer spins
are $1/2$. Hence their free energies are $9/2,13/2$ and their squared
norms are $1/4,3/4$.
Repeated faces give $W_p^2-1$ of energy $8$; edge-disjoint distinct faces
give $W_pW_q$ of energy $6$. All unordered pairs are thus retained in
\[
 B=-M/18+\sum_p\frac{W_p^2-1}{24}
  +\sum_{\partial p\cap\partial q=\varnothing}\frac{W_pW_q}{9}
  +\sum_{|\partial p\cap\partial q|=1}
       \left(\frac4{27}P_{pq,0}+\frac4{39}P_{pq,1}\right).
\]
Inserting this expression in \eqref{eq:fabel-local-D}, the first-order
electric term equals the first-order magnetic term because
$E_f(\mathcal W/3)=\sum_pf_pW_p$, where $E_f=\sum_ef_eE_e$.
Centering subtracts the expectation of the same vacuum.
Repeated and edge-disjoint face pairs cancel at order two.
For an adjacent pair with shared edge $e$, the weighted electric
eigenvalues are $3(f_p+f_q)-3f_e/2$ and $3(f_p+f_q)+f_e/2$.
The remaining second-order state is therefore
\begin{align}
 \Xi_f&=\kappa\xi^2 Z_f+O_{H^k,L,f}(\kappa\xi^3),\nonumber\\
 Z_f&=\sum_{\{p,q\}\ {\rm adjacent}}\delta_{pq}(f)
         (P_{pq,0}/9-P_{pq,1}/39),\quad
 \delta_{pq}(f)=f_p+f_q-2f_e. \label{eq:fabel-adjacent-state}
\end{align}
Different adjacent pairs have different six-edge outer boundaries.
Otherwise their symmetric difference would be a nonempty mod-two
two-cycle with at most four elementary faces. A face in that cycle
needs four distinct other faces to cancel its four edges (two distinct
square faces share at most one edge), a contradiction.
Central sign reversal of a boundary link present in only one pair
therefore kills the cross integral. It also kills the cross integral
with $H_0$ inserted. Both spin channels remain orthogonal. Thus
\begin{align}
 \|Z_f\|^2&=\frac{49}{13689}\mathcal D_f,&
 \langle Z_f,H_0Z_f\rangle&=\frac2{117}\mathcal D_f,&
 \mathcal D_f&=\sum_{\{p,q\}\ {\rm adjacent}}\delta_{pq}(f)^2 .
 \label{eq:fabel-D-moments}
\end{align}
The constants are respectively $1/(4\cdot9^2)+3/(4\cdot39^2)$ and
$(9/2)/(4\cdot9^2)+3(13/2)/(4\cdot39^2)$.
The central transformation
$U_i(n)\mapsto(-1)^{\sum_{j<i}n_j}U_i(n)$ negates every $W_p$ and
preserves the link Casimirs. It identifies the centered problems at
$\xi$ and $-\xi$, including the centered weighted state.
Consequently norm and energy are even:
\begin{equation}\label{eq:fabel-first-moments}
 \|\Xi_f\|^2=\frac{49}{13689}\kappa^2\xi^4\mathcal D_f+
 O_{L,f}(\kappa^2\xi^6),\quad
 q_{\mathcal A}[\Xi_f]=\frac2{117}\kappa^3\xi^4\mathcal D_f+
 O_{L,f}(\kappa^3\xi^6).
\end{equation}

Here are complete all-box counts for the new quadratic profiles.
For $e=(n,i)$, put $m=n+\mathbf e_i/2$.
A face in transverse direction $j$ on side $\epsilon$ has center
$m+\epsilon\mathbf e_j/2$. Averaging its four edge-midpoint quadratic
values gives
\[
 f_p=f_S(e)+\epsilon(Sm)_j+(S_{ii}+3S_{jj})/8.
\]
Indeed its midpoint displacements from its center are
$\pm\mathbf e_i/2,\pm\mathbf e_j/2$; the average correction to the
center value is $(S_{ii}+S_{jj})/8$.
Consequently opposite coplanar faces give
$\delta=(S_{ii}+3S_{jj})/4$.
Two hinged faces in directions $j,k$ on sides $\epsilon,\eta$ give
\[
 \delta=\epsilon(Sm)_j+\eta(Sm)_k+c_i,\qquad
 c_i=(2S_{ii}+3S_{jj}+3S_{kk})/8 .
\]
For each longitudinal coordinate the count and second midpoint moment
are $N_i=2L$ and $I_2=L(4L^2-1)/6$.
For each oriented transverse pair $(n,\epsilon)$ satisfying
$-L\le n,n+\epsilon\le L$, the sums are
\[
 N=4L,\quad \sum n=\sum\epsilon=0,\quad
 N_2=\sum n^2=\frac{2L(2L^2+1)}3,\quad
 E_n=\sum\epsilon n=-2L .
\]
These follow from the sums over $n=-L,\ldots,L-1$ at $\epsilon=1$
and $n=-L+1,\ldots,L$ at $\epsilon=-1$.
Expansion of every square now yields the exact finite formula
\begin{align*}
 \mathcal D_S^{\rm cop}&=N_i(2L-1)(2L+1)
                  \sum_{i\ne j}(S_{ii}+3S_{jj})^2/16,\\
 \mathcal D_S^{\rm hinge}
 &=\sum_i\bigl[
 I_2N^2(S_{ij}^2+S_{ik}^2)
 +N_iNN_2(S_{jj}^2+S_{kk}^2+2S_{jk}^2)\\
 &\hspace{12mm}+2N_iE_n^2(S_{jj}S_{kk}+S_{jk}^2)
 +2c_iN_iNE_n(S_{jj}+S_{kk})+c_i^2N_iN^2\bigr].
\end{align*}
In each summand $\{j,k\}=\{1,2,3\}\setminus\{i\}$ and
$\mathcal D_S=\mathcal D_S^{\rm cop}+\mathcal D_S^{\rm hinge}$.
There is no periodic boundary replacement.

Substitution of the three seed matrices gives
\begin{align}
 \mathcal D_0&=4L(2L-1)(2L+1)(4L^2+3),\nonumber\\
 \mathcal D_{\rm d}&=L(256L^4+74L^2-21)/6,\nonumber\\
 \mathcal D_{\rm o}&=64L^3(2L^2+1)/3.\label{eq:fabel-three-D}
\end{align}
All are positive for $L\ge2$.
The same substitution for $S_*=\eta\eta^{\mathsf T}$, or the already
proved cubic decomposition of the coefficient Gram matrix, gives
\begin{equation}\label{eq:fabel-rank-one-D}
 \mathcal D_{S_*}
 =\frac{L(14536832L^4+5453566L^2-1614327)}{24576}>0 .
\end{equation}
Equation \eqref{eq:fabel-first-moments} proves nonzero seed states for
sufficiently small positive $\xi$ on each fixed box and proves
\[
 \frac{q_{\mathcal A}[\Xi_{f_{S_*}}]}{\|\Xi_{f_{S_*}}\|^2}
       =\frac{234}{49}\kappa+O_L(\kappa\xi^2).
\]
This is an ordered fixed-box coefficient result, not a volume-uniform
remainder bound. In original variables
$\kappa^2\xi^4=1/(64g^{12}a^2)$ and
$\kappa^3\xi^4=1/(32g^{10}a^3)$.
The $L^5$ growth in \eqref{eq:fabel-rank-one-D} belongs to the computed
Taylor coefficient of both raw moments. Their common spatial and
material growth does not by itself reduce the quotient.

\subsection{Full Schur map and the remaining collective spectral question}
For a physical state $\Xi_f$ in the original Haar space set
$h_f=\Xi_f/\psi$, which is smooth at each finite regulator, and retain
\[
 F_f=J_m^*h_f,\quad
 v_f=\mathcal R_mF_f,\quad
 k_f=Q_mh_f+C^{-1}B_mF_f.
\]
These are the exact maps and domains of Theorem~9.1.
The identity $h_f=v_f+k_f$ has
\begin{align*}
 q(h_f,h_f)&=s_{m,0}(F_f,F_f)+c_m(k_f,k_f),\\
 \|\Xi_f\|^2&=\|v_f\|^2+\|k_f\|^2
                         +2\operatorname{Re}\langle v_f,k_f\rangle,\\
 \|v_f\|^2&=\|F_f\|^2+\int t^{-2}\,d\sigma_{F_f}(t),\\
 s_{m,0}(F_f,F_f)&=a_m(F_f,F_f)-\int t^{-1}\,d\sigma_{F_f}(t).
\end{align*}
No interaction or norm term has been removed.
The signed permutations above preserve the skeleton with $m\mid2L$:
the transverse coarse coordinates $-L+m\{0,\ldots,2L/m\}$ are invariant
under sign reversal, and chain products reverse by group inversion.
Vacuum conditional averaging therefore intertwines these symmetries by
change of variables in its defining integral. So do $B_m,C$ and the
full memory kernels; for the latter use functional calculus for
$e^{-sC}$ or $(C+z)^{-1}$. Their spectral matrix elements have the same
symmetry-forced vanishing cross terms. This does not evaluate the
remaining same-sector measures by a free approximation.

In the proved regime $\xi\le1/49152$, every nonzero constructed physical
state obeys $q/\|\Xi\|^2\ge\kappa(3-512\xi)\ge287\kappa/96$.
Outside that regime the exact three-measure formula remains valid but
does not compute the unknown seed measures. On any regulator sequence,
at least one of the nonzero three seed states has quotient no larger
than the quadratic combination, since \eqref{eq:fabel-sector-measure}
makes its quotient a weighted mean with weights $w_jd_j\ge0$.
Thus a vanishing quotient in this full-tensor quadratic family would
also give a vanishing quotient in a selected seed sector along that
sequence. At fixed regulator \eqref{eq:fabel-tensor-limit-quotient}
instead gives the precise endpoint measure and ratio. No exchange of
the material endpoint, infinite-volume limit, or spatial continuum
limit has been made. The calculation leaves those three interacting
seed-sector limits as substantive research, not as assumed estimates.
Section~\ref{sec:fabel-low} now calculates their complete fixed-box
weak-coupling pair measures, proves an injective lowest-band tensor map,
and constructs an actual positive-coupling projected Fabel sequence with
a vanishing energy quotient. It retains the whole-state raw norm and
band fraction. The common interacting continuum identification remains
unresolved.

\input{fabel_low_mode_addendum.tex}
\input{local_energy_band_addendum.tex}
\input{full_local_spectrum_addendum.tex}
\input{nonabelian_vertex_addendum.tex}
\input{interacting_tensor_band_addendum.tex}

\clearpage
\begin{thebibliography}{9}
\bibitem{DSS}
G.~Dusson, I.~M.~Sigal, and B.~Stamm,
\emph{The Feshbach--Schur map and perturbation theory},
in \emph{Partial Differential Equations, Spectral Theory, and Mathematical
Physics}, EMS Series of Congress Reports (2021), pp.~65--88.
\href{https://doi.org/10.4171/ECR/18-1/5}{doi:10.4171/ECR/18-1/5};
\href{https://arxiv.org/abs/2105.02058v1}{arXiv:2105.02058v1},
Theorem~1.2 and Eqs.~(1.10)--(1.16).
\bibitem{BL}
D.~Bakry and M.~Ledoux,
\emph{A logarithmic Sobolev form of the Li--Yau parabolic inequality},
Revista Matem\'atica Iberoamericana \textbf{22} (2006), no.~2, 683--702.
\href{https://doi.org/10.4171/RMI/470}{doi:10.4171/RMI/470},
Eq.~(1.7) and its proof, pp.~686--687.
\bibitem{Simon}
B.~Simon, \emph{A Feynman--Kac Formula for Unbounded Semigroups},
\href{https://arxiv.org/abs/math-ph/9907022v1}{arXiv:math-ph/9907022v1}
(1999), Theorem~1.1 and Eqs.~(1.1)--(1.3).
\bibitem{Mori}
H.~Mori, \emph{Transport, Collective Motion, and Brownian Motion},
Progress of Theoretical Physics \textbf{33} (1965), 423--455.
\href{https://doi.org/10.1143/PTP.33.423}{doi:10.1143/PTP.33.423}.
\bibitem{KS}
J.~Kogut and L.~Susskind,
\emph{Hamiltonian formulation of Wilson's lattice gauge theories},
Physical Review D \textbf{11} (1975), 395--408.
\href{https://doi.org/10.1103/PhysRevD.11.395}{doi:10.1103/PhysRevD.11.395}.
\bibitem{PW}
F.~Peter and H.~Weyl,
\emph{Die Vollst\"andigkeit der primitiven Darstellungen einer
geschlossenen kontinuierlichen Gruppe},
Mathematische Annalen \textbf{97} (1927), 737--755.
\href{https://doi.org/10.1007/BF01447892}{doi:10.1007/BF01447892}.
\bibitem{Yarotsky}
D.~A.~Yarotsky, \emph{Quasi-particles in weak perturbations of
non-interacting quantum lattice systems},
\href{https://arxiv.org/abs/math-ph/0411042v1}{arXiv:math-ph/0411042v1}
(2004), Section 2.
\bibitem{Baez}
J.~C.~Baez, \emph{Spin Network States in Gauge Theory},
Advances in Mathematics \textbf{117} (1996), 253--272.
\href{https://doi.org/10.1006/aima.1996.0012}{doi:10.1006/aima.1996.0012};
\href{https://arxiv.org/abs/gr-qc/9411007v1}{arXiv:gr-qc/9411007v1}.
\end{thebibliography}
\end{document}
